\cleardoublepage
\markboth{CHAPITRE IV}{CHAPITRE IV}
\oldpage[IV-1]{222}
\thispagestyle{plain}

\begin{center}
{\Large CHAPITRE IV}

\vspace{1.5cm}
{\LARGE\bfseries ÉTUDE LOCALE DES SCHÉMAS\\[0.3em]
ET DES MORPHISMES DE SCHÉMAS}

\vspace{1.5cm}
{\large\bfseries Sommaire}
\end{center}

\medskip
\begin{tabular}{@{}r p{0.84\textwidth}}
§ 1. & Conditions de finitude relatives. Ensembles constructibles dans les préschémas.\\
§ 2. & Changement de base et platitude.\\
§ 3. & Cycles premiers associés et décompositions primaires.\\
§ 4. & Changement du corps de base dans les préschémas.\\
§ 5. & Dimension et profondeur dans les préschémas.\\
§ 6. & Morphismes plats de préschémas localement noethériens.\\
§ 7. & Application aux relations entre un anneau local noethérien et son complété.
Anneaux excellents.\\
§ 8. & Limites projectives de préschémas.\\
§ 9. & Propriétés constructibles.\\
§ 10. & Préschémas de Jacobson.\\
§ 11.\footnotemark[1] & Propriétés topologiques des morphismes plats de présentation finie.
Critères locaux de platitude.\\
§ 12. & Étude des fibres des morphismes plats de présentation finie.\\
§ 13. & Morphismes équidimensionnels.\\
§ 14. & Morphismes universellement ouverts.\\
§ 15. & Étude des fibres d'un morphisme universellement ouvert.\\
§ 16. & Invariants différentiels. Morphismes différentiellement lisses.\\
§ 17. & Morphismes lisses, non ramifiés, étales.\\
§ 18. & Compléments sur les morphismes étales. Anneaux locaux henséliens.\\
§ 19. & Immersions régulières et transversalement régulières.\\
§ 20. & Sections hyperplanes~; projections génériques.\\
§ 21. & Prolongements infinitésimaux.\\
\end{tabular}
\footnotetext[1]{L'ordre et le contenu des §§ 11 à 21 ne sont donnés qu'à titre
indicatif et pourront être modifiés avant leur publication.}

\clearpage
\oldpage[IV-1]{223}

Les sujets traités dans ce Chapitre appellent les remarques suivantes~:

\medskip
\noindent
\emph{a)} Leur caractère commun est d'être relatifs à des propriétés
\emph{locales} de préschémas ou de morphismes, i.e. considérées en un point,
ou aux points d'une fibre, ou dans un voisinage (non spécifié) d'un point ou
d'une fibre. Ces propriétés sont généralement de nature \emph{topologique},
\emph{différentielle} ou \emph{dimensionnelle} (i.e. faisant intervenir les
notions de \emph{dimension} et de \emph{profondeur}) et sont liées aux
propriétés des \emph{anneaux locaux} aux points considérés. Un problème type
est de mettre en relation, pour un morphisme donné $f:X\to Y$ et un point
$x\in X$, les propriétés de $X$ en $x$ avec celles de $Y$ en $y=f(x)$ et de
la fibre $X_y=f^{-1}(y)$ en $x$. Un autre est de déterminer la nature
topologique (par exemple la constructibilité, ou le fait d'être ouvert ou
fermé) de l'ensemble des points $x\in X$ en lesquels $X$ possède une certaine
propriété ou pour lesquels la fibre $X_{f(x)}$ passant par $x$ possède une
certaine propriété en $x$. De même, on s'intéressera à la nature topologique
de l'ensemble des points $y\in Y$ tels que $X$ ait une certaine propriété en
tous les points de la fibre $X_y$, ou tels que cette fibre ait une certaine
propriété.

\medskip
\noindent
\emph{b)} Les notions les plus importantes pour les chapitres suivants sont
celle de \emph{morphisme plat de présentation finie} et celles de
\emph{morphisme lisse} et de \emph{morphisme étale}, qui en sont des cas
particuliers. Leur étude détaillée (jointe à celle de questions connexes)
commence à proprement parler au §~11.

\medskip
\noindent
\emph{c)} Les §§~1 à 10 peuvent être considérés comme de nature préliminaire
et développent trois types de techniques, utilisées non seulement dans les
autres paragraphes du chapitre mais aussi bien entendu dans les chapitres
suivants~:

\medskip
\noindent
\emph{c 1)} Aux §§~1 à 4 sont envisagés divers aspects de la notion de
\emph{changement de base}, surtout en relation avec des conditions de
\emph{finitude} ou de \emph{platitude}~; on y amorce la technique de
\emph{descente} sous ses aspects les plus élémentaires (les questions
d'«~effectivité~» liées à cette technique seront étudiées dans le chap.~V).

\medskip
\noindent
\emph{c 2)} Les §§~5 à 7 sont centrés sur ce que l'on peut appeler des
\emph{techniques noethériennes}, les préschémas considérés étant toujours
supposés localement noethériens, tandis qu'au contraire aucune condition de
finitude n'est généralement imposée aux \emph{morphismes}~; cela tient
essentiellement à ce que les notions de dimension et de profondeur ne sont
guère maniables que pour les anneaux locaux noethériens. Rappelons que le §~7
constitue une théorie «~fine~» des anneaux locaux noethériens, assez peu
utilisée dans la suite du chapitre.

\medskip
\noindent
\emph{c 3)} Les §§~8 à 10 donnent entre autres le moyen d'\emph{éliminer les
hypothèses noethériennes} sur les préschémas étudiés, en leur substituant des
hypothèses convenables de \emph{finitude} (la «~présentation finie~») sur les
\emph{morphismes} considérés~: l'avantage de cette substitution est que ces
dernières hypothèses sont \emph{stables par changement de base}, ce qui n'est
pas le cas pour les hypothèses noethériennes sur les préschémas. La technique
permettant ce passage repose, d'une part, sur l'utilisation de la notion de
\emph{limite projective} de préschémas, grâce à laquelle on peut ramener une
question à la même question sous des hypothèses \emph{noethériennes}~; d'autre
part, sur l'emploi systématique des \emph{ensembles constructibles}, qui ont
le double intérêt de se conserver par image réciproque (par un morphisme

\oldpage[IV-1]{224}

quelconque) et par image directe (par un morphisme de présentation finie), et
d'avoir des propriétés topologiques maniables dans les préschémas localement
noethériens. Les mêmes techniques permettent même souvent de se ramener à des
anneaux noethériens plus particuliers, par exemple les
$\mathbf Z$-\emph{algèbres de type fini}, et c'est là que les propriétés des
anneaux «~excellents~» (étudiées au §~7) interviennent de façon décisive.
Indépendamment d'ailleurs de la question de l'élimination des hypothèses
noethériennes, les techniques des §§~8 à 10, de nature très élémentaire, sont
d'usage constant dans presque toutes les applications.

\setcounter{section}{0}

\section{Conditions de finitude relatives. Ensembles constructibles dans les préschémas}
\label{section:IV.1-fr}

Nous allons, dans ce paragraphe, reprendre, en les complétant, l'exposé des
«~conditions de finitude~» pour un morphisme de préschémas $f:X\to Y$, donné
dans (I, 6.3 et 6.6). Il y a essentiellement deux notions de «~finitude~» de
nature \emph{globale} sur $X$, celle de morphisme \emph{quasi-compact}
(définie dans (I, 6.6.1)) et celle de morphisme \emph{quasi-séparé}~; il y a
d'autre part deux notions de «~finitude~» de nature \emph{locale} sur $X$,
celle de morphisme \emph{localement de type fini} (définie dans (I, 6.6.2))
et celle de morphisme \emph{localement de présentation finie}. En combinant
ces notions locales aux notions globales précédentes, on obtient les notions
de morphisme \emph{de type fini} (définie dans (I, 6.3.1)) et de morphisme
\emph{de présentation finie}. Pour la commodité du lecteur, nous donnerons de
nouveau, dans ce paragraphe, les principales propriétés énoncées dans (I, 6.3
et 6.6), en renvoyant bien entendu à ces numéros du chapitre I pour leurs
démonstrations.

Aux n\textsuperscript{os}~(1.8) et (1.9), nous complétons, dans le cadre des
préschémas, et en faisant usage des notions de finitude précédentes, les
résultats sur les ensembles constructibles donnés dans
(0\textsubscript{III}, §~9).

\subsection{Morphismes quasi-compacts}
\label{subsection:IV.1.1-fr}

\begin{definition}[1.1.1]
\label{IV.1.1.1-fr}
On dit qu'un morphisme de préschémas $f:X\to Y$ est \emph{quasi-compact} si
l'application continue $f$ de l'espace topologique $X$ dans l'espace
topologique $Y$ est quasi-compacte (0\textsubscript{III}, 9.1.1), autrement
dit si l'image réciproque $f^{-1}(U)$ de tout ouvert quasi-compact $U$ de $Y$
est quasi-compacte (cf. (I, 6.6.1)).
\end{definition}

Si $\mathfrak B$ est une base de la topologie de $Y$ formée d'ouverts
affines, pour que $f$ soit quasi-compact, il faut et il suffit que pour tout
$V\in\mathfrak B$, $f^{-1}(V)$ soit \emph{réunion finie d'ouverts affines}.
Par exemple, si $Y$ est affine et $X$ quasi-compact, \emph{tout} morphisme
$f:X\to Y$ est quasi-compact (I, 6.6.1).

Si $f:X\to Y$ est un morphisme quasi-compact, il est clair que pour tout
ouvert $V$ de $Y$, la restriction de $f$ à $f^{-1}(V)$ est un morphisme
quasi-compact $f^{-1}(V)\to V$. Inversement, si $(U_\alpha)$ est un
recouvrement ouvert de $Y$ et $f:X\to Y$ un morphisme tel que les restrictions
$f^{-1}(U_\alpha)\to U_\alpha$ soient quasi-compactes, alors $f$ est
quasi-compact. Par

\oldpage[IV-1]{225}

suite, si $f:X\to Y$ est un $S$-morphisme de $S$-préschémas, et s'il existe
un recouvrement ouvert $(S_\lambda)$ de $S$ tel que les restrictions
$g^{-1}(S_\lambda)\to h^{-1}(S_\lambda)$ de $f$ (où $g$ et $h$ sont les
morphismes structuraux) soient quasi-compactes, alors $f$ est quasi-compact.

\begin{proposition}[1.1.2]
\label{IV.1.1.2-fr}
\begin{enumerate}
  \item[(i)] Une immersion $X\to Y$ est quasi-compacte si elle est fermée, ou si
  l'espace sous-jacent à $Y$ est localement noethérien ou si l'espace
  sous-jacent à $X$ est noethérien.
  \item[(ii)] Le composé de deux morphismes quasi-compacts est quasi-compact.
  \item[(iii)] Si $f:X\to Y$ est un $S$-morphisme quasi-compact, il en est de même de
  $f_{(S')}:X_{(S')}\to Y_{(S')}$ pour toute extension $g:S'\to S$ du
  préschéma de base.
  \item[(iv)] Si $f:X\to X'$ et $g:Y\to Y'$ sont deux $S$-morphismes
  quasi-compacts,
  \[
    f\times_S g:X\times_S Y\to X'\times_S Y'
  \]
  est quasi-compact.
  \item[(v)] Si le composé $g\circ f$ de deux morphismes $f:X\to Y$, $g:Y\to Z$
  est quasi-compact, et si $g$ est séparé, ou l'espace sous-jacent à $X$
  localement noethérien, $f$ est quasi-compact.
  \item[(vi)] Pour qu'un morphisme $f$ soit quasi-compact, il faut et il suffit que
  $f_{\mathrm{red}}$ le soit.
\end{enumerate}
\end{proposition}

Pour la démonstration, voir (I, 6.6.4). On notera que l'assertion (vi) est
aussi conséquence de la proposition plus générale~:

\begin{proposition}[1.1.3]
\label{IV.1.1.3-fr}
Soient $f:X\to Y$, $g:Y\to Z$ deux morphismes. Si $g\circ f$ est
quasi-compact et $f$ surjectif, $g$ est quasi-compact.
\end{proposition}

\begin{proof}
En effet, si $V$ est un ouvert quasi-compact dans $Z$,
$f^{-1}(g^{-1}(V))$ est quasi-compact par hypothèse, et l'on a
$g^{-1}(V)=f(f^{-1}(g^{-1}(V)))$, puisque $f$ est surjectif~; donc
$g^{-1}(V)$ est quasi-compact.
\end{proof}

\begin{corollary}[1.1.4]
\label{IV.1.1.4-fr}
Soient $f:X\to Y$, $g:Y'\to Y$ deux morphismes, $X'=X\times_Y Y'$,
$f'=f_{(Y')}:X'\to Y'$~; supposons $g$ quasi-compact et surjectif (resp.
surjectif). Alors, pour que $f$ soit quasi-compact (resp. dominant), il suffit
que $f'$ le soit.
\end{corollary}

\begin{proof}
Si $g':X'\to X$ est la projection canonique, $g'$ est un morphisme surjectif
(I, 3.5.2), et l'on a $f\circ g'=g\circ f'$~; si $f'$ est quasi-compact
(resp. dominant), il en est de même de $g\circ f'$ puisque $g$ est
quasi-compact (1.1.2) (resp. surjectif)~; comme $f\circ g'$ est quasi-compact
(resp. dominant) et que $g'$ est surjectif, on en déduit que $f$ est
quasi-compact (1.1.3) (resp. dominant).
\end{proof}

\oldpage[IV-1]{226}

Pour abréger, nous appellerons \emph{point maximal} d'un préschéma $X$ tout
point générique d'une composante irréductible de $X$~; si
$X=\operatorname{Spec}(A)$ est affine, les points maximaux de $X$ sont les
\emph{idéaux premiers minimaux} de $A$.

\begin{proposition}[1.1.5]
\label{IV.1.1.5-fr}
Soit $f:X\to Y$ un morphisme quasi-compact. Les conditions suivantes sont
équivalentes~:
\begin{enumerate}
  \item[a)] $f$ est dominant~;
  \item[b)] pour tout point maximal $y$ de $Y$, on a
  $f^{-1}(y)\neq\varnothing$~;
  \item[c)] pour tout point maximal $y$ de $Y$, $f^{-1}(y)$ contient un point
  maximal de $X$.
\end{enumerate}
\end{proposition}

L'équivalence de \emph{a)} et \emph{c)} a été démontrée dans (I, 6.6.5). Il
est clair que \emph{c)} entraîne \emph{b)}~; d'autre part, si $X'$ est une
composante irréductible de $X$ rencontrant $f^{-1}(y)$, $f(X')$ est
irréductible et contient $y$, donc est contenu dans la composante irréductible
$Y'=\overline{\{y\}}$ de $Y$ (0\textsubscript{I}, 2.1.6)~; si $x$ est le
point générique de $X'$, on a donc nécessairement $f(x)=y$
(0\textsubscript{I}, 2.1.5)~; donc \emph{b)} entraîne \emph{c)}.

\begin{proposition}[1.1.6]
\label{IV.1.1.6-fr}
Soient $f':X'\to Y$, $f'':X''\to Y$ deux morphismes, $X$ le préschéma somme
$X'\coprod X''$, $f:X\to Y$ le morphisme égal à $f'$ dans $X'$ et à $f''$
dans $X''$. Pour que $f$ soit quasi-compact, il faut et il suffit que $f'$ et
$f''$ le soient.
\end{proposition}

Cela résulte aussitôt de la définition.

\subsection{Morphismes quasi-séparés}
\label{subsection:IV.1.2-fr}

\begin{definition}[1.2.1]
\label{IV.1.2.1-fr}
Soient $X$, $Y$ deux préschémas. On dit qu'un morphisme $f:X\to Y$ est
\emph{quasi-séparé} (ou que $X$ est \emph{quasi-séparé sur} $Y$) si le
morphisme diagonal
$\Delta_f=(1_X,1_X)_Y:X\to X\times_Y X$ est quasi-compact. On dit qu'un
préschéma $X$ est \emph{quasi-séparé} s'il est quasi-séparé sur
$\operatorname{Spec}(\mathbf Z)$.
\end{definition}

Par définition, tout morphisme séparé est quasi-séparé, $\Delta_f$ étant une
immersion fermée (1.1.2, (i))~; un \emph{schéma} $X$ est un préschéma
quasi-séparé, étant séparé sur $\operatorname{Spec}(\mathbf Z)$ (I, 5.4.1).

\begin{proposition}[1.2.2]
\label{IV.1.2.2-fr}
\begin{enumerate}
  \item[(i)] Tout monomorphisme de préschémas $f:X\to Y$ (en particulier toute
  immersion) est quasi-séparé.
  \item[(ii)] Si $f:X\to Y$ et $g:Y\to Z$ sont deux morphismes quasi-séparés,
  $g\circ f:X\to Z$ est quasi-séparé.
  \item[(iii)] Soient $X$, $Y$ deux $S$-préschémas, $f:X\to Y$ un $S$-morphisme
  quasi-séparé. Alors, pour tout changement de base $g:S'\to S$, le morphisme
  $f_{(S')}:X_{(S')}\to Y_{(S')}$ est quasi-séparé.
  \item[(iv)] Si $f:X\to Y$, $f':X'\to Y'$ sont deux $S$-morphismes quasi-séparés,
  alors
  \[
    f\times_S f':X\times_S X'\to Y\times_S Y'
  \]
  est quasi-séparé.
  \item[(v)] Si $f:X\to Y$, $g:Y\to Z$ sont deux morphismes tels que $g\circ f$
  soit quasi-séparé, alors $f$ est quasi-séparé.
  \item[(vi)] Si $f:X\to Y$ est un morphisme quasi-séparé,
  $f_{\mathrm{red}}:X_{\mathrm{red}}\to Y_{\mathrm{red}}$ est quasi-séparé.
\end{enumerate}
\end{proposition}

\begin{proof}
En vertu de (I, 5.5.12), il suffit de démontrer (i), (ii) et (iii).

L'assertion (i) est triviale, tout monomorphisme étant séparé (I, 5.5.1).
Pour démontrer (iii), on se ramène aussitôt au cas où $Y=S$ (I, 3.3.11) et
l'assertion résulte de ce que
$\Delta_{f_{(S')}}=(\Delta_f)_{(S')}$ (I, 5.3.4) et de (1.1.2, (iii)). Pour
démontrer (ii), considérons les projections $p$ et $q$ de $X\times_Y X$ sur
$X$~; si $i=(p,q)_Z$, on sait que le diagramme
\[
\xymatrix{
  X\times_Y X \ar[r]^-{i} \ar[d]_-{\pi} & X\times_Z X
    \ar[d]^-{f\times_Z f} \\
  Y \ar[r]_-{\Delta_g} & Y\times_Z Y
}
\]
\oldpage[IV-1]{227}
(où $\pi$ est le morphisme structural) est commutatif et identifie
$X\times_Y X$ au produit des $(Y\times_Z Y)$-préschémas $Y$ et
$X\times_Z X$ (I, 5.3.5). Si $g$ est quasi-séparé, $\Delta_g$ est
quasi-compact, donc il en est de même de $i$ (1.1.2, (iii))~; si de plus $f$
est quasi-séparé, $\Delta_f$ est quasi-compact, donc il en est de même de
$i\circ\Delta_f$ (1.1.2, (ii)) qui est égal à $\Delta_{g\circ f}$.
\end{proof}

\begin{corollary}[1.2.3]
\label{IV.1.2.3-fr}
\begin{enumerate}
  \item[(i)] Soit $X$ un préschéma quasi-séparé. Alors tout morphisme $f:X\to Y$
  est quasi-séparé.
  \item[(ii)] Soit $Y$ un préschéma quasi-séparé. Pour qu'un morphisme $f:X\to Y$
  soit quasi-séparé, il faut et il suffit que le préschéma $X$ soit
  quasi-séparé.
\end{enumerate}
\end{corollary}

Cela résulte aussitôt de (1.2.2, (ii) et (v)) appliqué à $X$, $Y$ et
$\operatorname{Spec}(\mathbf Z)$.

\begin{proposition}[1.2.4]
\label{IV.1.2.4-fr}
Soient $f:X\to Y$ un morphisme, $g:Y\to Z$ un morphisme quasi-séparé. Si
$g\circ f$ est quasi-compact, il en est de même de $f$.
\end{proposition}

\begin{proof}
On sait que $f$ est le morphisme composé
$X\xrightarrow{\Gamma_f}X\times_ZY\xrightarrow{\mathrm{pr}_2}Y$
(I, 5.3.13)~; d'autre part, $\mathrm{pr}_2$ s'identifie à
$(g\circ f)\times_Z1_Y$ (I, 3.3.4) et si $g\circ f$ est quasi-compact, il en
est de même de $\mathrm{pr}_2$ (1.1.2, (iv)). Enfin, on a le diagramme
commutatif
\[
\xymatrix{
  X \ar[r]^-{\Gamma_f} \ar[d]_-{f} & X\times_ZY
    \ar[d]^-{f\times 1_Y} \\
  Y \ar[r]_-{\Delta_g} & Y\times_ZY
}
\tag{1.2.4.1}
\label{IV.1.2.4.1-fr}
\]
qui identifie $X$ au produit des $(Y\times_ZY)$-préschémas $Y$ et
$X\times_ZY$ (I, 5.3.7). Comme par hypothèse $\Delta_g$ est un morphisme
quasi-compact, il en est de même de $\Gamma_f$ (1.1.2, (iii))~; la conclusion
résulte donc de (1.1.2, (ii)).
\end{proof}

\begin{proposition}[1.2.5]
\label{IV.1.2.5-fr}
Soient $f:X\to Y$, $g:Y'\to Y$ deux morphismes,
$X'=X\times_Y Y'$, $f'=f_{(Y')}:X'\to Y'$. Si $g$ est quasi-compact et
surjectif et $f'$ quasi-séparé, alors $f$ est quasi-séparé.
\end{proposition}

\begin{proof}
On a en effet $(X'\times_{Y'}X')=(X\times_YX)_{(Y')}$ et
$\Delta_{f'}=(\Delta_f)_{(Y')}$ (I, 5.3.4)~; comme la projection
$X'\times_{Y'}X'\to X\times_YX$ est quasi-compacte (1.1.2, (iii)) et
surjective (I, 3.5.2), on peut, en vertu de (I, 3.3.11), appliquer (1.1.4),
qui montre que puisque $\Delta_{f'}$ est un morphisme quasi-compact, il en est
de même de $\Delta_f$.
\end{proof}

\begin{proposition}[1.2.6]
\label{IV.1.2.6-fr}
Soient $f:X\to Y$ un morphisme, $(U_\alpha)$ un recouvrement de $Y$ par des
ouverts tels que les préschémas induits $U_\alpha$ soient quasi-séparés. Pour
que $f$ soit quasi-séparé, il faut et il suffit que chacun des préschémas
induits par $X$ sur les ouverts $f^{-1}(U_\alpha)$ soit quasi-séparé.
\end{proposition}

\begin{proof}
L'image réciproque dans $X\times_YX$ de $U_\alpha$ est
$X_\alpha\times_{U_\alpha}X_\alpha$, en posant
$X_\alpha=f^{-1}(U_\alpha)$, et la restriction
$X_\alpha\to X_\alpha\times_{U_\alpha}X_\alpha$ de $\Delta_f$ n'est autre
que $\Delta_{f_\alpha}$, en désignant par $f_\alpha$ la restriction
$X_\alpha\to U_\alpha$ de $f$~; en vertu de la définition (1.2.1) et du
caractère local de la notion de morphisme quasi-compact (1.1.1), pour que $f$
soit quasi-séparé, il faut et il suffit que chacun des morphismes $f_\alpha$
le soit. Mais comme par hypothèse le morphisme
$U_\alpha\to\operatorname{Spec}(\mathbf Z)$ est quasi-séparé, il revient au
même de dire que $f_\alpha$ est quasi-séparé ou que le composé
$X_\alpha\xrightarrow{f_\alpha}U_\alpha\to\operatorname{Spec}(\mathbf Z)$
l'est (1.2.2, (ii) et (v)), d'où la proposition.
\end{proof}

\oldpage[IV-1]{228}

Pour vérifier qu'un morphisme est quasi-séparé, on est donc ramené à donner
des critères pour qu'un préschéma soit quasi-séparé~:

\begin{proposition}[1.2.7]
\label{IV.1.2.7-fr}
Soient $X$ un préschéma, $(U_\alpha)$ un recouvrement de $X$ formé d'ouverts
quasi-compacts. Les propriétés suivantes sont équivalentes~:
\begin{enumerate}
  \item[a)] $X$ est quasi-séparé.
  \item[b)] Pour tout ouvert quasi-compact $U$ de $X$, l'immersion canonique
  $U\to X$ est quasi-compacte (autrement dit, $U$ est rétrocompact dans $X$).
  \item[b')] L'intersection de deux ouverts quasi-compacts de $X$ est
  quasi-compacte.
  \item[c)] Pour tout couple d'indices $(\alpha,\beta)$, l'intersection
  $U_\alpha\cap U_\beta$ est quasi-compacte.
\end{enumerate}
\end{proposition}

\begin{proof}
Les propriétés \emph{b)} et \emph{b')} sont trivialement équivalentes par
définition d'un morphisme quasi-compact. Comme un ouvert quasi-compact est
réunion finie d'ouverts affines, pour deux ouverts quasi-compacts $U$, $V$ de
$X$, $U\times_{\mathbf Z}V$ est un ouvert quasi-compact de
$X\times_{\mathbf Z}X$ (I, 3.2.7), dont l'image réciproque par $\Delta_X$ est
$U\cap V$, donc \emph{a)} entraîne \emph{b')}. Il est trivial que \emph{b')}
entraîne \emph{c)}~; enfin, si \emph{c)} est vérifiée, les
$U_\alpha\times_{\mathbf Z}U_\beta$ forment un recouvrement de
$X\times_{\mathbf Z}X$ par des ouverts quasi-compacts, et l'on sait que pour
que le morphisme $\Delta_X:X\to X\times_{\mathbf Z}X$ soit quasi-compact, il
suffit que les images réciproques $U_\alpha\cap U_\beta$ des
$U_\alpha\times_{\mathbf Z}U_\beta$ par $\Delta_X$ soient quasi-compactes
(1.1.1)~; donc \emph{c)} entraîne \emph{a)}.
\end{proof}

\begin{corollary}[1.2.8]
\label{IV.1.2.8-fr}
Tout préschéma $X$ dont l'espace sous-jacent est localement noethérien (par
exemple un préschéma localement noethérien) est quasi-séparé~; tout morphisme
$f:X\to Y$ est alors quasi-séparé.
\end{corollary}

\begin{proposition}[1.2.9]
\label{IV.1.2.9-fr}
Soient $X'$, $X''$ deux préschémas, $X=X'\coprod X''$ leur somme,
$f':X'\to Y$, $f'':X''\to Y$ deux morphismes, $f:X\to Y$ le morphisme qui
coïncide avec $f'$ dans $X'$ et avec $f''$ dans $X''$. Pour que $f$ soit
quasi-séparé il faut et il suffit que $f'$ et $f''$ le soient.
\end{proposition}

\begin{proof}
En effet, $X\times_YX$ est somme des quatre préschémas
$X'\times_YX'$, $X'\times_YX''$, $X''\times_YX'$ et $X''\times_YX''$, et
$\Delta_f$ est le morphisme qui coïncide avec $\Delta_{f'}$ dans $X'$ et avec
$\Delta_{f''}$ dans $X''$~; la proposition résulte donc aussitôt des
définitions.
\end{proof}

\subsection{Morphismes localement de type fini}
\label{subsection:IV.1.3-fr}

\begin{env}[1.3.1]
\label{IV.1.3.1-fr}
Rappelons que si $A$ est un anneau, une $A$-algèbre (commutative) $B$ est dite
\emph{de type fini} si elle est engendrée par un nombre fini d'éléments de
$B$, ou, ce qui revient au même, si elle est isomorphe au quotient d'une
algèbre de polynômes $A[T_1,\ldots,T_n]$ par un idéal de cette algèbre. Il est
clair que pour toute $A$-algèbre (commutative) $A'$,
$B\otimes_AA'$ est alors une $A'$-algèbre de type fini. Si $B$ est une
$A$-algèbre de type fini et $C$ une $B$-algèbre de type fini, $C$ est une
$A$-algèbre de type fini, car si $C$ est un quotient de
$B[T_1,\ldots,T_n]$ et $B$ un quotient de $A[S_1,\ldots,S_m]$,
$B[T_1,\ldots,T_n]=B\otimes_AA[T_1,\ldots,T_n]$ est un quotient de
$A[S_1,\ldots,S_m,T_1,\ldots,T_n]$, donc il en est de même de $C$.
\end{env}

\begin{definition}[1.3.2]
\label{IV.1.3.2-fr}
Soient $f:X\to Y$ un morphisme de préschémas, $x$ un point de $X$,
$y=f(x)$. On dit que $f$ est \emph{de type fini en $x$} s'il existe un
voisinage ouvert affine $V$ de $y$ et un voisinage ouvert affine $U$ de $x$
tels que $f(U)\subset V$ et que l'anneau $A(U)$ soit une $A(V)$-algèbre de
type fini. On dit que $f$ est \emph{localement de type fini} s'il est de type
fini en tout point de $X$.
\end{definition}

\oldpage[IV-1]{229}

\begin{proposition}[1.3.3]
\label{IV.1.3.3-fr}
Si $Y$ est un préschéma localement noethérien et $f:X\to Y$ un morphisme
localement de type fini, alors $X$ est localement noethérien.
\end{proposition}

Pour la démonstration, voir (I, 6.3.7).

\begin{proposition}[1.3.4]
\label{IV.1.3.4-fr}
\begin{enumerate}
  \item[(i)] Toute immersion locale est localement de type fini.
  \item[(ii)] Si deux morphismes $f:X\to Y$, $g:Y\to Z$ sont localement de
  type fini, il en est de même de $g\circ f$.
  \item[(iii)] Si $f:X\to Y$ est un $S$-morphisme localement de type fini,
  $f_{(S')}:X_{(S')}\to Y_{(S')}$ est localement de type fini pour toute
  extension $S'\to S$ du préschéma de base.
  \item[(iv)] Si $f:X\to X'$ et $g:Y\to Y'$ sont deux $S$-morphismes
  localement de type fini, $f\times_Sg$ est localement de type fini.
  \item[(v)] Si le composé $g\circ f$ de deux morphismes est localement de
  type fini, $f$ est localement de type fini.
  \item[(vi)] Si un morphisme $f$ est localement de type fini, il en est de
  même de $f_{\mathrm{red}}$.
\end{enumerate}
\end{proposition}

Pour la démonstration, voir (I, 6.6.6).

\begin{corollary}[1.3.5]
\label{IV.1.3.5-fr}
Soit $f:X\to Y$ un morphisme localement de type fini. Pour tout morphisme
$Y'\to Y$ tel que $Y'$ soit localement noethérien, $X\times_YY'$ est
localement noethérien.
\end{corollary}

\begin{proof}
En effet, $f_{(Y')}:X\times_YY'\to Y'$ est localement de type fini par
(1.3.4, (iii)), et il suffit d'appliquer (1.3.3).
\end{proof}

\begin{proposition}[1.3.6]
\label{IV.1.3.6-fr}
Soit $\varphi:A\to B$ un homomorphisme d'anneaux. Pour que le morphisme
correspondant $f:\operatorname{Spec}(B)\to\operatorname{Spec}(A)$ soit localement de type fini, il faut et
il suffit que $B$ soit une $A$-algèbre de type fini.
\end{proposition}

Pour la démonstration, voir (I, 6.3.3), en tenant compte de ce que $f$ est
quasi-compact et de (I, 6.6.3).

\begin{proposition}[1.3.7]
\label{IV.1.3.7-fr}
Pour qu'un morphisme localement de type fini $f:X\to Y$ soit surjectif, il
faut et il suffit que, pour tout corps algébriquement clos $\Omega$,
l'application $X(\Omega)\to Y(\Omega)$ correspondant à $f$ (I, 3.4.1) soit
surjective.
\end{proposition}

Pour la démonstration, voir (I, 6.3.10).

\begin{env}[1.3.8]
\label{IV.1.3.8-fr}
Soit $A$ un anneau. On dit qu'une $A$-algèbre $B$ est \emph{essentiellement de
type fini} si $B$ est $A$-isomorphe à une $A$-algèbre de la forme $S^{-1}C$,
où $C$ est une $A$-algèbre \emph{de type fini} et $S$ une partie multiplicative
de $C$.
\end{env}

\par\medskip\noindent
\begin{proposition}[1.3.9]
\label{IV.1.3.9-fr}
\begin{enumerate}
  \item[(i)] Si $B$ est une $A$-algèbre essentiellement de type fini et $C$
  une $B$-algèbre essentiellement de type fini, alors $C$ est une
  $A$-algèbre essentiellement de type fini.
  \item[(ii)] Soient $B$ une $A$-algèbre essentiellement de type fini, $A'$
  une $A$-algèbre~; alors $B'=B\otimes_AA'$ est une $A'$-algèbre
  essentiellement de type fini.
\end{enumerate}
\end{proposition}

\begin{proof}
\begin{enumerate}
  \item[(i)] Soient $B=S'^{-1}B'$ et $C=T'^{-1}C'$, où $B'$ (resp. $C'$) est
  une $A$-algèbre (resp. une $B$-algèbre) de type fini et $S'$ (resp. $T'$)
  une partie multiplicative de $B'$ (resp. $C'$)~; alors $C'$ est de la forme
  $S'^{-1}C''$, où $C''$ est une $B'$-algèbre de type fini, donc $C$ est aussi
  un anneau de fractions de $C''$~; comme $C''$ est une $A$-algèbre de type
  fini, cela prouve notre assertion.
  \item[(ii)] Si $B$ est anneau de fractions d'une $A$-algèbre de type fini
  $C$, $B'$ est un anneau de fractions de la $A'$-algèbre de type fini
  $C\otimes_AA'$, d'où (ii).
\end{enumerate}
\end{proof}

\oldpage[IV-1]{230}

\begin{env}[1.3.10]
\label{IV.1.3.10-fr}
Si $B$ est une $A$-algèbre \emph{locale} essentiellement de type fini, $B$ est
de la forme $C_{\mathfrak q}$, où $C$ est une $A$-algèbre de type fini et
$\mathfrak q$ un idéal premier de $C$ (Bourbaki, \emph{Alg. comm.}, chap. II,
\S~5, prop. 11). Soit $\mathfrak p$ l'image réciproque dans $A$ de l'idéal
$\mathfrak q$~; si l'on pose $S=A-\mathfrak p$, $C_{\mathfrak q}$ est aussi un
anneau local en un idéal premier de $S^{-1}C$~; comme $S^{-1}C$ est une
algèbre de type fini sur $A_{\mathfrak p}=S^{-1}A$, on voit que $B$ est aussi
une $A_{\mathfrak p}$-algèbre \emph{essentiellement de type fini}, et en outre
l'homomorphisme $A_{\mathfrak p}\to B$ est \emph{local}.
\end{env}

\begin{proposition}[1.3.11]
\label{IV.1.3.11-fr}
Si $B$ est une $A$-algèbre locale essentiellement de type fini, $B$ est
$A$-isomorphe à une $A$-algèbre quotient d'une $A$-algèbre de la forme
$C_{\mathfrak q}$, où $C$ est une algèbre de polynômes
$A[T_1,\ldots,T_n]$ et $\mathfrak q$ un idéal premier de $C$.
\end{proposition}

\begin{proof}
En effet, par définition, $B$ est isomorphe à $C'_{\mathfrak q'}$, où $C'$ est
une $A$-algèbre de type fini et $\mathfrak q'$ un idéal premier de $C'$~; mais
$C'=C/\mathfrak b$, où $C=A[T_1,\ldots,T_n]$ et $\mathfrak b$ est un idéal de
$C$~; donc $\mathfrak q'=\mathfrak q/\mathfrak b$, où $\mathfrak q$ est un
idéal premier de $C$, et $C'_{\mathfrak q'}$ est isomorphe à
$C_{\mathfrak q}/\mathfrak bC_{\mathfrak q}$.
\end{proof}

\subsection{Morphismes localement de présentation finie}
\label{subsection:IV.1.4-fr}

\begin{env}[1.4.1]
\label{IV.1.4.1-fr}
Étant donné un anneau $A$, nous dirons qu'une $A$-algèbre $B$ (commutative)
est \emph{de présentation finie} si elle est isomorphe au quotient
$A[T_1,\ldots,T_n]/\mathfrak a$ d'une algèbre de polynômes sur $A$ par un
idéal \emph{de type fini} $\mathfrak a$ de $A[T_1,\ldots,T_n]$. Pour toute
$A$-algèbre (commutative) $A'$, $B\otimes_AA'$ est alors une $A'$-algèbre de
présentation finie, étant isomorphe au quotient de $A'[T_1,\ldots,T_n]$ par
l'image canonique dans cet anneau de $\mathfrak a\otimes_AA'$, qui est
évidemment un $A'[T_1,\ldots,T_n]$-module de type fini. Si $B$ est une
$A$-algèbre de présentation finie et $C$ une $B$-algèbre de présentation
finie, alors $C$ est une $A$-algèbre de présentation finie. En effet, soient
$B=A[S_1,\ldots,S_m]/\mathfrak a$ et
$C=B[T_1,\ldots,T_n]/\mathfrak b$, où $\mathfrak a$ (resp. $\mathfrak b$) est
un idéal de type fini de $A[S_1,\ldots,S_m]$ (resp.
$B[T_1,\ldots,T_n]$)~; l'anneau $B[T_1,\ldots,T_n]$ est isomorphe à
$A[S_1,\ldots,S_m,T_1,\ldots,T_n]/\mathfrak a'$, où $\mathfrak a'$ est
l'image canonique de
$\mathfrak a\otimes_AA[T_1,\ldots,T_n]$, et est donc un idéal de type fini de
$A[S_1,\ldots,S_m,T_1,\ldots,T_n]$~; l'idéal $\mathfrak b$ de
$B[T_1,\ldots,T_n]$ est de la forme $\mathfrak b'/\mathfrak a'$, où
$\mathfrak b'$ est un idéal de $A[S_1,\ldots,S_m,T_1,\ldots,T_n]$~; comme
$\mathfrak a'$ et $\mathfrak b'/\mathfrak a'$ sont des modules de type fini
sur $A[S_1,\ldots,S_m,T_1,\ldots,T_n]$, il en est de même de $\mathfrak b'$,
et $C$, isomorphe à
$A[S_1,\ldots,S_m,T_1,\ldots,T_n]/\mathfrak b'$, est donc bien une
$A$-algèbre de présentation finie. Notons enfin que si $A$ est
\emph{noethérien}, il en est de même de $A[T_1,\ldots,T_n]$, donc toute
$A$-algèbre de type fini est alors de présentation finie.
\end{env}

\begin{definition}[1.4.2]
\label{IV.1.4.2-fr}
Soient $f:X\to Y$ un morphisme de préschémas, $x$ un point de $X$, $y=f(x)$.
On dit que $f$ est \emph{de présentation finie en $x$} s'il existe un voisinage
ouvert affine $V$ de $y$ et un voisinage ouvert affine $U$ de $x$ tels que
$f(U)\subset V$ et que l'anneau $A(U)$ soit une $A(V)$-algèbre de présentation
finie. On dit que $f$ est \emph{localement de présentation finie} s'il est de
présentation finie en tout point de $X$.
\end{definition}

Si $Y$ est localement noethérien, il revient au même de dire que $f$ est
localement de type fini ou localement de présentation finie.

\begin{proposition}[1.4.3]
\label{IV.1.4.3-fr}
\begin{enumerate}
  \item[(i)] Tout isomorphisme local est localement de présentation finie.
  \item[(ii)] Si deux morphismes $f:X\to Y$, $g:Y\to Z$ sont localement de
  présentation finie, il en est de même de $g\circ f$.

\oldpage[IV-1]{231}

  \item[(iii)] Si $f:X\to Y$ est un $S$-morphisme localement de présentation
  finie, $f_{(S')}:X_{(S')}\to Y_{(S')}$ est localement de présentation finie
  pour toute extension $S'\to S$ du préschéma de base.
  \item[(iv)] Si $f:X\to X'$ et $g:Y\to Y'$ sont deux $S$-morphismes
  localement de présentation finie, $f\times_Sg$ est localement de
  présentation finie.
  \item[(v)] Si le composé $g\circ f$ de deux morphismes est localement de
  présentation finie et si $g$ est localement de type fini, alors $f$ est
  localement de présentation finie.
\end{enumerate}
\end{proposition}

\begin{proof}
L'assertion (i) est triviale. Pour démontrer (iii), on peut se limiter au cas
où $S=Y$ (I, 3.3.11)~; soient $x'$ un point de $X_{(S')}$, $s'$ et $x$ ses
projections sur $S'$ et $X$ respectivement, $s$ la projection commune de $s'$
et $x$ sur $S$. Par hypothèse, il y a un voisinage ouvert affine $V$ (resp.
$U$) de $s$ (resp. $x$) tels que l'image de $U$ soit contenue dans $V$ et que
$A(U)$ soit une $A(V)$-algèbre de présentation finie~; soit $V'$ un voisinage
ouvert affine de $s'$ dont l'image dans $S$ soit contenue dans $V$~; alors
$U\times_VV'$ est un voisinage ouvert affine de $x'$ dont l'image dans $S'$
est contenue dans $V'$ et dont l'anneau $A(U)\otimes_{A(V)}A(V')$ est une
$A(V')$-algèbre de présentation finie (1.4.1).

Pour démontrer (ii), considérons un point $x\in X$, et soient $y=f(x)$,
$z=g(f(x))$. Il y a par hypothèse un voisinage ouvert affine $W=\operatorname{Spec}(C)$
(resp. $V=\operatorname{Spec}(B)$) de $z$ (resp. $y$) tels que $g(V)\subset W$ et que $B$
soit une $C$-algèbre de présentation finie. D'autre part, il y a un voisinage
ouvert affine $V'=\operatorname{Spec}(B')$ (resp. $U=\operatorname{Spec}(A)$) de $y$ (resp. $x$) tels que
$f(U)\subset V'$ et que $A$ soit une $B'$-algèbre de présentation finie. Soit
$s\in B$ tel que l'ouvert $D(s)=\operatorname{Spec}(B_s)$ dans $V$ soit un voisinage de $y$
contenu dans $V'$, et soit $U'=f^{-1}(D(s))\subset U$~; on a
$U'=\operatorname{Spec}(A\otimes_{B'}B_s)$, et $A\otimes_{B'}B_s$ est une $B_s$-algèbre de
présentation finie (1.4.1)~; d'aute part $B_s=B[T]/(1-sT)$ est une
$B$-algèbre de présentation finie par définition~; par suite (1.4.1),
$A\otimes_{B'}B_s$ est une $C$-algèbre de présentation finie, ce qui prouve
(ii).

On sait que (iv) résulte de (i), (ii) et (iii) (I, 3.5.1). Pour prouver (v),
considérons le diagramme commutatif
\[
\xymatrix{
X \ar[r]^{\Gamma_f} \ar[d] & X\times_ZY \ar[d]^{f\times 1} \\
Y \ar[r]_{\Delta_g} & Y\times_ZY
}
\]
qui identifie $X$ au produit des $(Y\times_ZY)$-préschémas $Y$ et
$X\times_ZY$ (I, 5.3.7, $\Delta_g$ étant le morphisme diagonal). Si l'on sait
que $\Delta_g$ est localement de présentation finie, il en résultera par
(iii) qu'il en est de même de $\Gamma_f$. D'autre part, on a la factorisation
de $f$ en
$X\xrightarrow{\Gamma_f}X\times_ZY\xrightarrow{p_2}Y$ (I, 5.3.13), où la
projection $p_2$ est égale à $(g\circ f)\times_Z1_Y$~; comme $g\circ f$ est
supposé localement de présentation finie, il en est de même de $p_2$ par
(iv), et il résultera finalement de (ii) que $f$ est aussi localement de
présentation finie. On voit donc qu'on est ramené à démontrer le
\end{proof}

\begin{corollary}[1.4.3.1]
\label{IV.1.4.3.1-fr}
Soit $g:Y\to Z$ un morphisme localement de type fini~; alors le morphisme
diagonal $\Delta_g:Y\to Y\times_ZY$ est localement de présentation finie.
\end{corollary}

\oldpage[IV-1]{232}

\begin{proof}
On peut se borner au cas où $Z=\operatorname{Spec}(A)$, $Y=\operatorname{Spec}(B)$, $B$ étant une
$A$-algèbre de type fini~; le morphisme $\Delta_g$ correspond alors à
l'\emph{homomorphisme d'augmentation} $\delta:B\otimes_AB\to B$ tel que
$\delta(x\otimes y)=xy$. Compte tenu de la définition (1.4.2), il suffit de
montrer que le noyau $\mathfrak J$ de $\delta$ est un idéal \emph{de type
fini}, ce qui résulte de (0, 20.4.4).
\end{proof}

\begin{proposition}[1.4.4]
\label{IV.1.4.4-fr}
Soient $A$ un anneau, $B$ une $A$-algèbre, $B'$ une algèbre de polynômes
$A[T_1,\ldots,T_n]$, $u:B'\to B$ un homomorphisme surjectif de
$A$-algèbres. Pour que $B$ soit une $A$-algèbre de présentation finie, il faut
et il suffit que le noyau $\mathfrak J$ de $u$ soit un idéal de type fini de
$B'$.
\end{proposition}

\begin{proof}
La condition est suffisante par définition (1.4.1). Pour voir qu'elle est
nécessaire, remarquons que le morphisme $g:\operatorname{Spec}(B')\to\operatorname{Spec}(A)$ est
localement de type fini~; si $B$ est une $A$-algèbre de présentation finie,
le morphisme $f:\operatorname{Spec}(B)\to\operatorname{Spec}(B')$ correspondant à $u$ est tel que
$g\circ f:\operatorname{Spec}(B)\to\operatorname{Spec}(A)$ soit localement de présentation finie, donc il
résulte de (1.4.3, (v)) que $f$ est aussi localement de présentation finie.
Or, pour démontrer que l'idéal $\mathfrak J$ est de type fini, il suffit de
prouver que $\widetilde{\mathfrak J}$ est un $\widetilde{B'}$-Module de type
fini (II, 6.1.4.1). Revenant à la définition (1.4.2), on est finalement ramené
à prouver le
\end{proof}

\begin{lemma}[1.4.4.1]
\label{IV.1.4.4.1-fr}
Soient $\mathfrak a$ un idéal d'un anneau $C$, $s$ un élément de $C$, tel que
$C_s/\mathfrak a_s$ soit une $C$-algèbre de présentation finie~; alors
$\mathfrak a_s$ est un idéal de type fini dans l'anneau $C_s$.
\end{lemma}

\begin{proof}
Par hypothèse, il existe une $C$-algèbre de polynômes
$C'=C[T_1,\ldots,T_n]$ et un $C$-homomorphisme surjectif
$v:C'\to C_s/\mathfrak a_s$ dont le noyau $\mathfrak b$ est de type fini.
Soit $p:C_s\to C_s/\mathfrak a_s$ l'homomorphisme canonique~; pour tout $i$,
il y a un $t_i\in C$ tel que $p(t_i/s^k)=v(T_i)$ (on peut supposer que
l'exposant $k$ est le même pour tous les indices $i$). Considérons alors la
$C_s$-algèbre $D=C_s[T_1,\ldots,T_n]$, et l'homomorphisme de $C_s$-algèbres
$w:D\to C_s$ tel que $w(T_i)=t_i/s^k$~; $w$ est évidemment surjectif, et il
en est par suite de même de $v'=p\circ w:D\to C_s/\mathfrak a_s$. D'autre
part, tout polynôme de $D$ peut s'écrire $P/s^m$, où
$P\in C'=C[T_1,\ldots,T_n]$, et l'on a
$v'(P/s^m)=(1/s^m)v(P)$~; comme $1/s$ est inversible dans $C_s$, la relation
$v'(P/s^m)=0$ dans $C_s/\mathfrak a_s$ équivaut à $v(P)=0$, et par suite le
noyau $\mathfrak b'$ de $v'$ est engendré par l'image canonique de
$\mathfrak b$ dans $D$, et \emph{a fortiori} est de type fini dans l'anneau
$D$. Mais
$\mathfrak b'=v'^{-1}(0)=w^{-1}(p^{-1}(0))=w^{-1}(\mathfrak a_s)$, et comme
$w$ est surjectif, $\mathfrak a_s=w(\mathfrak b')$, ce qui prouve que
$\mathfrak a_s$ est un idéal de type fini de $C_s$.
\end{proof}

\begin{corollary}[1.4.5]
\label{IV.1.4.5-fr}
Soient $X$, $Y$ deux préschémas, $j:X\to Y$ une immersion, $U$ un ouvert de
$Y$ tel que $j(X)$ soit fermé dans $U$, $\mathcal J$ l'idéal quasi-cohérent de
$\mathcal O_U$ qui définit le sous-préschéma fermé de $Y$ associé à $j$. Pour
que $j$ soit localement de présentation finie, il faut et il suffit que
$\mathcal J$ soit un $\mathcal O_U$-Module de type fini.
\end{corollary}

\par\medskip\noindent
\begin{proposition}[1.4.6]
\label{IV.1.4.6-fr}
Soit $\varphi:A\to B$ un homomorphisme d'anneaux. Pour que le morphisme
correspondant $f:\operatorname{Spec}(B)\to\operatorname{Spec}(A)$ soit localement de présentation finie,
il faut et il suffit que $B$ soit une $A$-algèbre de présentation finie.
\end{proposition}

\begin{proof}
La condition étant trivialement suffisante, il reste à prouver qu'elle est
nécessaire. Si $f$ est localement de présentation finie, il résulte déjà de
(1.3.6) que $B$ est une $A$-algèbre de type fini, donc il existe un
homomorphisme surjectif $u:B'=A[T_1,\ldots,T_n]\to B$ de $A$-algèbres~; il
résulte de (1.4.3, (v)) que l'immersion
$j:\operatorname{Spec}(B)\to\operatorname{Spec}(B')$ est localement

\oldpage[IV-1]{233}

de présentation finie~; donc (1.4.5), si $\mathfrak J$ est le noyau de $u$,
le $\widetilde{B'}$-Module $\widetilde{\mathfrak J}$ est de type fini, et par
suite (II, 6.1.4.1) $\mathfrak J$ est un idéal de type fini.
\end{proof}

\begin{proposition}[1.4.7]
\label{IV.1.4.7-fr}
Soit $\rho:A\to B$ un homomorphisme d'anneaux faisant de $B$ une $A$-algèbre
finie. Pour que $B$ soit une $A$-algèbre de présentation finie, il faut et il
suffit que $B$ soit un $A$-module de présentation finie.
\end{proposition}

Nous utiliserons le

\begin{lemma}[1.4.7.1]
\label{IV.1.4.7.1-fr}
Si $B$ est une $A$-algèbre finie, il existe un $A$-homomorphisme surjectif
d'algèbres $u:B'\to B$, où $B'$ est une $A$-algèbre finie et de présentation
finie qui est un $A$-module libre.
\end{lemma}

\begin{proof}
En effet, il y a un système fini $(a_i)_{1\leq i\leq m}$ de générateurs de la
$A$-algèbre $B$ tel que chaque $a_i$ vérifie une relation $F_i(a_i)=0$, où
$F_i\in A[X]$ est un polynôme unitaire de degré $>0$~; la $A$-algèbre quotient
$B'_i=A[X]/(F_i)$ est donc libre de rang fini sur $A$~; soit $c_i$ la classe de
$X$ dans $B'_i$. Il existe un $A$-homomorphisme d'algèbres $u_i:B'_i\to B$
tel que $u_i(c_i)=a_i$~; il suffit alors de prendre pour $B'$ le produit
tensoriel $B'_1\otimes_AB'_2\otimes_A\cdots\otimes_AB'_m$ et de considérer
l'homomorphisme $u_1\otimes u_2\otimes\cdots\otimes u_m:B'\to B$, qui est
surjectif par construction. En outre, si $B''=A[T_1,\ldots,T_m]$ et si $b'_i$
désigne l'image canonique de $c_i$ dans $B'$, le $A$-homomorphisme d'algèbres
$v:B''\to B'$ tel que $v(T_i)=b'_i$ ($1\leq i\leq m$) est surjectif, et son
noyau $\mathfrak b'$ est l'idéal de $B''$ engendré par les $F_i(T_i)$, donc de
type fini~; ceci prouve que $B'$ est une $A$-algèbre de présentation finie.
\end{proof}

\begin{proof}
Ce lemme étant démontré, gardons les mêmes notations, et soit
$w=v\circ u:B''\to B$. Si $\mathfrak b$ (resp. $\mathfrak a$) est le noyau de
$w$ (resp. $u$), on a $\mathfrak a=v(\mathfrak b)$ puisque $v$ est surjectif.
Or (1.4.4), si $B$ est une $A$-algèbre de présentation finie, $\mathfrak b$
est un idéal de type fini de $B''$, donc $\mathfrak a$ est un idéal de type
fini de $B'$, donc un $A$-module de type fini puisque $B'$ est une $A$-algèbre
finie~; comme $B'$ est un $A$-module libre, $B$ est un $A$-module de
présentation finie. Inversement, si $B'$ est un $A$-module de présentation
finie, $\mathfrak a$ est un $A$-module de type fini (Bourbaki,
\emph{Alg. comm.}, chap. I, \S~2, n\textsuperscript{o}~8, lemme 9) et
\emph{a fortiori} un idéal de type fini de $B'$~; par suite, $B$ est par
définition une $B'$-algèbre de présentation finie, et comme $B'$ est une
$A$-algèbre de présentation finie, $B$ est une $A$-algèbre de présentation
finie.
\end{proof}

\subsection{Morphismes de type fini}
\label{subsection:IV.1.5-fr}

\begin{proposition}[1.5.1]
\label{IV.1.5.1-fr}
Soient $f:X\to Y$ un morphisme, $(U_\alpha)$ un recouvrement de $Y$ formé
d'ouverts affines. Les conditions suivantes sont équivalentes~:
\begin{enumerate}
  \item[a)] $f$ est localement de type fini et quasi-compact.
  \item[b)] Pour tout $\alpha$, $f^{-1}(U_\alpha)$ est réunion finie d'ouverts
  affines $V_{\alpha i}$ tels que chaque anneau $A(V_{\alpha i})$ soit une
  $A(U_\alpha)$-algèbre de type fini.
  \item[c)] Pour tout ouvert affine $U$ de $Y$, $f^{-1}(U)$ est réunion finie
  d'ouverts affines $V_j$ tels que les anneaux $A(V_j)$ soient des
  $A(U)$-algèbres de type fini.
\end{enumerate}
\end{proposition}

Pour la démonstration, voir (I, 6.3.2 et 6.6.3).

\begin{definition}[1.5.2]
\label{IV.1.5.2-fr}
On dit qu'un morphisme $f$ est \emph{de type fini} s'il vérifie les conditions
équivalentes de la proposition (1.5.1).
\end{definition}

\oldpage[IV-1]{234}

\begin{proposition}[1.5.3]
\label{IV.1.5.3-fr}
Soit $f:X\to Y$ un morphisme de type fini~; si $Y$ est noethérien, il en est
de même de $X$.
\end{proposition}

Pour la démonstration, voir (I, 6.3.7).

\begin{proposition}[1.5.4]
\label{IV.1.5.4-fr}
\begin{enumerate}
  \item[(i)] Soit $j:X\to Y$ un morphisme d'immersion. Si $j$ est
  quasi-compact (ce qui est le cas si $j$ est une immersion fermée, ou si
  l'espace sous-jacent à $X$ est noethérien, ou si l'espace sous-jacent à $Y$
  est localement noethérien), $j$ est de type fini.
  \item[(ii)] Le composé de deux morphismes de type fini est de type fini.
  \item[(iii)] Si $f:X\to Y$ est un $S$-morphisme de type fini,
  $f_{(S')}:X_{(S')}\to Y_{(S')}$ est de type fini pour toute extension
  $S'\to S$ du préschéma de base.
  \item[(iv)] Si $f:X\to X'$ et $g:Y\to Y'$ sont deux $S$-morphismes de type
  fini, $f\times_Sg$ est de type fini.
  \item[(v)] Si le composé $g\circ f$ de deux morphismes est de type fini et
  si $g$ est quasi-séparé ou si $X$ est noethérien, $f$ est de type fini.
  \item[(vi)] Si un morphisme $f$ est de type fini, il en est de même de
  $f_{\mathrm{red}}$.
\end{enumerate}
\end{proposition}

Cela résulte aussitôt (sauf le cas (v) où $X$ est noethérien, démontré dans
(I, 6.3.6)) des définitions et de (1.1.2), (1.2.4) et (1.3.4).

\begin{corollary}[1.5.5]
\label{IV.1.5.5-fr}
Soit $f:X\to Y$ un morphisme de type fini. Pour tout morphisme $Y'\to Y$ tel
que $Y'$ soit noethérien, $X\times_YY'$ est noethérien.
\end{corollary}

Cela résulte de (1.5.4, (iii)) et de (1.5.3).

\begin{corollary}[1.5.6]
\label{IV.1.5.6-fr}
Soit $X$ un préschéma de type fini sur un préschéma localement noethérien $S$.
Alors tout $S$-morphisme $f:X\to Y$ est de type fini.
\end{corollary}

Pour la démonstration, voir (I, 6.3.9).

\begin{proposition}[1.5.7]
\label{IV.1.5.7-fr}
Soit $\varphi:A\to B$ un homomorphisme d'anneaux. Pour que le morphisme
correspondant $f:\operatorname{Spec}(B)\to\operatorname{Spec}(A)$ soit de type fini, il faut et il suffit
que $\varphi$ fasse de $B$ une $A$-algèbre de type fini.
\end{proposition}

Pour la démonstration, voir (I, 6.3.3).

\subsection{Morphismes de présentation finie}
\label{subsection:IV.1.6-fr}

\begin{definition}[1.6.1]
\label{IV.1.6.1-fr}
Soient $X$, $Y$ deux préschémas. On dit qu'un morphisme $f:X\to Y$ est
\emph{de présentation finie} s'il vérifie les trois conditions suivantes~:
\begin{enumerate}
  \item[(i)] $f$ est localement de présentation finie~;
  \item[(ii)] $f$ est quasi-compact (ce qui, lorsque (i) est vérifiée,
  équivaut à dire que $f$ est de type fini (1.5.2))~;
  \item[(iii)] $f$ est quasi-séparé.
\end{enumerate}
On dit alors aussi que $X$ est \emph{de présentation finie sur $Y$}, ou un
\emph{$Y$-préschéma de présentation finie}.
\end{definition}

La condition (iii) est automatiquement vérifiée si $f$ est séparé, ou si $X$
est localement noethérien (1.2.8)~; lorsque $Y$ est localement noethérien, il
revient au même de dire que $f$ est \emph{de présentation finie} ou que $f$
est \emph{de type fini} (cette dernière condition entraînant que $X$ est
localement noethérien (1.3.3)).

\oldpage[IV-1]{235}

\begin{proposition}[1.6.2]
\label{IV.1.6.2-fr}
\begin{enumerate}
  \item[(i)] Toute immersion quasi-compacte et localement de présentation
  finie (en particulier toute immersion ouverte quasi-compacte) est de
  présentation finie.
  \item[(ii)] Le composé de deux morphismes de présentation finie est de
  présentation finie.
  \item[(iii)] Soient $X$, $Y$ deux $S$-préschémas, $f:X\to Y$ un
  $S$-morphisme de présentation finie. Alors, pour tout morphisme $S'\to S$,
  le morphisme $f_{(S')}:X_{(S')}\to Y_{(S')}$ est de présentation finie.
  \item[(iv)] Soient $f:X\to Y$, $f':X'\to Y'$ deux $S$-morphismes de
  présentation finie~; alors
  $f\times_Sf':X\times_SX'\to Y\times_SY'$ est de présentation finie.
  \item[(v)] Si le composé $g\circ f$ de deux morphismes est de présentation
  finie et si $g$ est quasi-séparé et localement de présentation finie, $f$
  est de présentation finie.
\end{enumerate}
\end{proposition}

Cela résulte immédiatement de (1.1.2), (1.2.2), (1.2.4) et (1.4.3).

Il résulte en particulier de (1.6.2, (iii)) que si $f$ est un morphisme de
présentation finie et $U$ un ouvert de $Y$, la restriction
$f^{-1}(U)\to U$ de $f$ est un morphisme de présentation finie. Inversement,
soit $(U_\alpha)$ un recouvrement de $Y$ par des ouverts affines, et supposons
que les restrictions $f^{-1}(U_\alpha)\to U_\alpha$ de $f$ soient des
morphismes de présentation finie~; il en résulte alors que $f$ est de
présentation finie, car $f$ est évidemment localement de présentation finie,
quasi-compact en vertu de (1.1.1) et quasi-séparé en vertu de (1.2.6). On peut
donc dire que la propriété pour un morphisme $f:X\to Y$ d'être de présentation
finie est \emph{locale sur $Y$}.

Si $X$ est un préschéma quasi-séparé, tout morphisme $f:X\to Y$ est
quasi-séparé (1.2.3). Donc, si $f$ est quasi-compact et localement de
présentation finie, $f$ est alors de présentation finie.

En particulier~:

\begin{corollary}[1.6.3]
\label{IV.1.6.3-fr}
Soit $\varphi:A\to B$ un homomorphisme d'anneaux. Pour que le morphisme
correspondant $f:\operatorname{Spec}(B)\to\operatorname{Spec}(A)$ soit de présentation finie, il faut et
il suffit que $\varphi$ fasse de $B$ une $A$-algèbre de présentation finie.
\end{corollary}

La nécessité de la condition résulte en effet de (1.4.6).

\begin{remark}[1.6.4]
\label{IV.1.6.4-fr}
Dans la définition (1.6.1), la condition (iii) n'est pas conséquence des deux
autres. Soit par exemple $Y$ un schéma affine dont l'espace sous-jacent ne
soit pas noethérien, et soit $U$ un ouvert non quasi-compact dans $Y$. Soit
$X$ le préschéma obtenu en recollant deux préschémas $Y_1$, $Y_2$ isomorphes à
$Y$ le long des ouverts $U_1$, $U_2$ correspondant à $U$ (0$_{\mathrm I}$,
4.1.7), de sorte que $X$ est réunion de deux ouverts isomorphes respectivement
à $Y_1$ et $Y_2$ (et que nous identifions à ces derniers), avec
$Y_1\cap Y_2=U$. En outre, il y a un morphisme $f:X\to Y$ qui coïncide dans
$Y_i$ avec l'isomorphisme donné $Y_i\to Y$ ($i=1,2$). Il est clair que $f$
vérifie la condition (i) de (1.6.1), étant un isomorphisme local~; il vérifie
aussi (ii), l'image réciproque d'un ouvert quasi-compact de $Y$ étant la
réunion de ses images dans $Y_1$ et $Y_2$~; mais comme $Y_1\cap Y_2$ n'est pas
quasi-compact, $f$ n'est pas quasi-séparé (1.2.7).
\end{remark}

\begin{proposition}[1.6.5]
\label{IV.1.6.5-fr}
Soient $X'$, $X''$ deux préschémas, $X=X'\amalg X''$ leur somme,
$f':X'\to Y$, $f'':X''\to Y$ deux morphismes, $f:X\to Y$ le morphisme qui
coïncide avec $f'$ dans $X'$ et avec $f''$ dans $X''$. Pour que $f$ soit de
présentation finie, il faut et il suffit que $f'$ et $f''$ le soient.
\end{proposition}

\par\medskip\noindent
\begin{proof}
Il suffit de montrer que pour que $f$ possède une des trois propriétés de la
définition

\oldpage[IV-1]{236}

(1.6.1), il faut et il suffit que $f'$ et $f''$ la possèdent~; c'est évident
pour la propriété d'être localement de présentation finie, qui est locale sur
$X$~; pour la propriété d'être quasi-compact, cela a été vu en (1.1.6) et pour
la propriété d'être quasi-séparé, en (1.2.9).
\end{proof}

\subsection{Améliorations de résultats antérieurs}
\label{subsection:IV.1.7-fr}

Nous allons donner dans ce numéro une liste de propositions démontrées dans
les chapitres précédents, dont l'énoncé peut être amélioré à l'aide des
nouvelles conditions de finitude introduites ci-dessus.

\begin{env}[1.7.1]
\label{IV.1.7.1-fr}
Dans les énoncés de (I, 6.4.2, 6.4.3 et 6.4.9), on peut, au lieu de supposer
que $X$ et $Y$ sont de type fini sur le corps $K$, supposer seulement qu'ils
sont \emph{localement de type fini} sur $K$~; pour (I, 6.4.2), il suffit
d'observer que tout point d'un préschéma $X$ qui est fermé dans tout ouvert
affine de $X$ est fermé dans $X$, et un schéma affine localement de type fini
sur $K$ est \emph{ipso facto} de type fini (1.5.1). Pour (I, 6.4.9), on utilise
(1.3.7) et (1.3.4, (v)).
\end{env}

\begin{env}[1.7.2]
\label{IV.1.7.2-fr}
Dans (I, 6.5.1, (ii)), on peut au lieu de supposer que $S$ est localement
noethérien et $Y$ de type fini sur $S$, supposer seulement que $Y$ est
\emph{localement de présentation finie} sur $S$, comme le montre aussitôt la
démonstration (en appliquant la définition (1.4.2)). De même dans (I, 6.5.4,
(ii)) et (I, 6.5.5, (ii)) il suffit de supposer que $f$ est un $S$-morphisme,
$Y$ étant localement de présentation finie sur $S$.
\end{env}

\begin{env}[1.7.3]
\label{IV.1.7.3-fr}
Dans (I, 7.1.11), pour la deuxième assertion, au lieu de supposer $S$
localement noethérien et $Y$ de type fini sur $S$, on peut seulement supposer
$Y$ localement de présentation finie sur $S$ (la démonstration dépendant de
(I, 6.5.1, (ii)))~; de même dans (I, 7.1.12 à 7.1.15).
\end{env}

\begin{env}[1.7.4]
\label{IV.1.7.4-fr}
Dans (I, 9.2.2), on peut remplacer les trois hypothèses par la seule hypothèse
(entraînée par chacune des autres) que $f$ est \emph{quasi-compact et
quasi-séparé}, en vertu de (1.2.6) et (1.2.7). On notera que dans (I, 9.2.1),
l'hypothèse signifie exactement que $f$ est quasi-compact et quasi-séparé.
\end{env}

\begin{env}[1.7.5]
\label{IV.1.7.5-fr}
Dans (I, 9.3.1, 9.3.2 et 9.3.3), on peut remplacer les deux hypothèses sur $X$
par l'hypothèse (moins forte) que $X$ est un préschéma \emph{quasi-compact et
quasi-séparé}, la démonstration utilisant exactement ces deux propriétés (en
vertu de (1.2.7)).
\end{env}

\begin{env}[1.7.6]
\label{IV.1.7.6-fr}
Dans (I, 9.3.4 et 9.3.5), il suffit de supposer $X$ quasi-compact,
$\mathcal F$ et $\mathcal J$ quasi-cohérents et de type fini.
\end{env}

\begin{env}[1.7.7]
\label{IV.1.7.7-fr}
Dans (I, 9.4.7), on peut affaiblir l'hypothèse b) en supposant seulement $X$
\emph{quasi-séparé}, puisqu'il suffit seulement que l'immersion canonique
$U\to X$ soit quasi-compacte (1.2.7). On peut faire la même modification dans
(I, 9.4.8, 9.4.9 et 9.4.10).
\end{env}

\begin{env}[1.7.8]
\label{IV.1.7.8-fr}
Dans (I, 9.5.1), les conditions indiquées entre parenthèses dans l'énoncé
peuvent être remplacées par la condition que $f$ soit quasi-compact et
quasi-séparé.
\end{env}

\begin{env}[1.7.9]
\label{IV.1.7.9-fr}
Dans (I, 9.6.5), on peut remplacer les hypothèses sur $X$ par l'hypothèse
moins forte que $X$ est quasi-compact et quasi-séparé, comme le montre la
démonstration,

\oldpage[IV-1]{237}

qui n'utilise que (I, 9.4.7). Dans (I, 9.6.6), l'hypothèse «~$X$ est un
schéma quasi-compact~» peut être remplacée par «~$X$ est un préschéma
quasi-compact et quasi-séparé~».
\end{env}

\begin{env}[1.7.10]
\label{IV.1.7.10-fr}
Dans (II, 1.7.15), on peut supposer seulement que $Y$ est quasi-compact et
quasi-séparé, la démonstration n'utilisant que (I, 9.6.5).
\end{env}

\begin{env}[1.7.11]
\label{IV.1.7.11-fr}
Dans la seconde assertion de (II, 3.4.5), on peut substituer aux deux
hypothèses sur $Y$ l'hypothèse plus faible que $Y$ est quasi-compact et
quasi-séparé. De même dans (II, 3.4.8).
\end{env}

\begin{env}[1.7.12]
\label{IV.1.7.12-fr}
Dans (II, 3.8.5), on peut de même supposer seulement que $Y$ est quasi-compact
et quasi-séparé, cette hypothèse intervenant par (I, 9.4.7).
\end{env}

\begin{env}[1.7.13]
\label{IV.1.7.13-fr}
Dans (II, 4.4.1, 4.4.6 et 4.4.7), il suffit de supposer que $Y$ est
\emph{quasi-compact et quasi-séparé}, compte tenu de (1.2.3) dans la
démonstration de (II, 4.4.7). De même, dans (II, 4.4.10.1), il suffit de
supposer $Z$ quasi-compact et quasi-séparé.
\end{env}

\begin{env}[1.7.14]
\label{IV.1.7.14-fr}
Dans (II, 4.5.2) et (II, 4.5.5), il suffit de supposer $X$
\emph{quasi-compact et quasi-séparé}, cette hypothèse intervenant par (I,
9.3.1 et 9.3.2).
\end{env}

\begin{env}[1.7.15]
\label{IV.1.7.15-fr}
Dans (II, 4.6.8), il suffit encore de supposer $X$ \emph{quasi-compact et
quasi-séparé}, cette hypothèse intervenant par (I, 9.4.7).
\end{env}

\begin{env}[1.7.16]
\label{IV.1.7.16-fr}
Dans (II, 5.1.2), il suffit de supposer que $X$ est quasi-compact et
quasi-séparé (l'hypothèse intervenant par (II, 4.5.2 et 4.5.5)). De même dans
(II, 5.1.9), où on utilise (I, 9.6.5).
\end{env}

\begin{env}[1.7.17]
\label{IV.1.7.17-fr}
Dans (II, 5.2.1), il suffit toujours de supposer $X$ quasi-compact et
quasi-séparé (l'hypothèse intervenant par (II, 4.5.2)).
\end{env}

\begin{env}[1.7.18]
\label{IV.1.7.18-fr}
Dans (II, 5.2.2), il faut supposer $f$ quasi-compact et $X$ quasi-séparé~; la
démonstration actuelle est en effet insuffisante (avec les hypothèses faites),
car du fait que pour tout $\mathcal O_X$-Module quasi-cohérent $\mathcal F$ on
a $\mathrm R^1f_*(\mathcal F)=0$, il ne s'ensuit pas nécessairement, pour un
ouvert affine $U$ de $Y$, que la même propriété soit valable pour les
$(\mathcal O_X|f^{-1}(U))$-Modules quasi-cohérents, si on ne sait pas qu'un tel
Module est la \emph{restriction} d'un $\mathcal O_X$-Module quasi-cohérent (la
même remarque s'appliquant aux conditions b) et c'))~; lorsque $f$ est
quasi-compact et $X$ quasi-séparé, ce prolongement est possible en vertu de
(I, 9.4.7), modifié plus haut (1.7.7).
\end{env}

\begin{env}[1.7.19]
\label{IV.1.7.19-fr}
Dans (II, 5.3.2), il suffit de supposer $Y$ \emph{quasi-compact et
quasi-séparé} (ce qui intervient par (II, 4.4.6 et 4.4.7)). De même dans (II,
5.5.3, (ii)).
\end{env}

\begin{env}[1.7.20]
\label{IV.1.7.20-fr}
Dans (III, 1.2.2, 1.2.3 et 1.2.4), on peut remplacer les deux hypothèses sur
$X$ par l'hypothèse moins forte que $X$ est un préschéma \emph{quasi-compact
et quasi-séparé}, la démonstration n'utilisant que (I, 9.3.3).
\end{env}

\begin{env}[1.7.21]
\label{IV.1.7.21-fr}
Dans (III, 1.4.10 à 1.4.14), on peut remplacer «~séparé~» par
«~quasi-séparé~», cette hypothèse servant uniquement à pouvoir appliquer (I,
9.2.2) (voir ci-dessus (1.7.4)). De même, dans (III, 1.4.15), il suffit de
supposer que le morphisme $u$ est \emph{quasi-séparé et quasi-compact}.
\end{env}

\begin{env}[1.7.22]
\label{IV.1.7.22-fr}
Dans (III, 2.1.3), on peut encore remplacer les hypothèses par l'hypothèse
moins forte que $X$ est un préschéma quasi-compact et quasi-séparé, le
raisonnement étant le même que dans (III, 1.2.2).
\end{env}

\subsection{Morphismes de présentation finie et ensembles constructibles}
\label{subsection:IV.1.8-fr}

\oldpage[IV-1]{238}

\begin{env}[1.8.1]
\label{IV.1.8.1-fr}
Soit $X$ un préschéma~; on sait que tout ouvert quasi-compact dans $X$ est
réunion finie d'ouverts affines, et réciproquement. Pour qu'un ouvert $U$ de
$X$ soit \emph{rétrocompact} dans $X$ (0\textsubscript{III}, 9.1.1), il faut
et il suffit donc que pour tout \emph{ouvert affine} $V$ de $X$, $U\cap V$
soit quasi-compact. Lorsque $X$ est \emph{quasi-séparé} (1.2.1), tout ouvert
quasi-compact dans $X$ est rétrocompact dans $X$ (1.2.7), donc toute partie
localement constructible de $X$ est rétrocompacte dans $X$
(0\textsubscript{III}, 9.1.13)~; si \emph{de plus} $X$ est
\emph{quasi-compact}, il y a identité entre parties constructibles et parties
localement constructibles dans $X$ (0\textsubscript{III}, 9.1.12), et entre
parties ouvertes constructibles et parties ouvertes quasi-compactes
(0\textsubscript{III}, 9.1.5).
\end{env}

\begin{proposition}[1.8.2]
\label{IV.1.8.2-fr}
Soient $X$, $Y$ deux préschémas, $f:X\to Y$ un morphisme. Pour toute partie
constructible (resp. localement constructible) $Z$ de $Y$, $f^{-1}(Z)$ est
une partie constructible (resp. localement constructible) de $X$.
\end{proposition}

\begin{proof}
Supposons $Z$ localement constructible~; pour tout $x\in X$, il y aura un
voisinage ouvert affine $V$ de $f(x)$ tel que $Z\cap V$ soit constructible
dans $V$~; si la proposition est démontrée pour les parties constructibles,
$f^{-1}(Z)\cap f^{-1}(V)$ sera constructible dans $f^{-1}(V)$, donc
$f^{-1}(Z)$ sera localement constructible. Il suffit par suite de considérer
le cas où $Z$ est constructible, et compte tenu de (0\textsubscript{III},
9.1.3), on est ramené au cas où $Z$ est ouvert et \emph{rétrocompact} dans
$Y$, autrement dit tel que l'injection canonique $Z\to Y$ soit un morphisme
quasi-compact~; il suffit alors de voir que $f^{-1}(Z)$ est
\emph{rétrocompact dans $X$}, autrement dit tel que l'injection canonique
$f^{-1}(Z)\to X$ soit un morphisme quasi-compact~; mais cela résulte de
(1.1.2, (iii)), puisque $f^{-1}(Z)=Z\times_YX$ (I, 4.4.1).
\end{proof}

\begin{lemma}[1.8.3]
\label{IV.1.8.3-fr}
Soient $X$ un préschéma quasi-compact et quasi-séparé, $Z$ une partie
constructible de $X$. Il existe alors un schéma affine $X'$ et un morphisme
de présentation finie $f:X'\to X$ tels que $f(X')=Z$.
\end{lemma}

\begin{proof}
Comme $X$ est quasi-compact, il est réunion finie d'ouverts affines $X_i$
($i\in I$), et comme $X$ est quasi-séparé, il résulte de (1.2.7) que chacune
des immersions canoniques $g_i:X_i\to X$ est de présentation finie~; on en
conclut que si $f_i:X'_i\to X_i$ est un morphisme de présentation finie, il
en est de même de $g_i\circ f_i:X'_i\to X$ (1.6.2), et si $X'$ est le
préschéma somme des $X'_i$, le morphisme $h:X'\to X$ qui coïncide avec
$g_i\circ f_i$ dans chaque $X'_i$ est lui aussi de présentation finie
(1.6.5). Comme $Z\cap X_i$ est constructible dans $X_i$
(0\textsubscript{III}, 9.1.8), on voit qu'on est ramené à prouver le lemme
lorsque $X_i$ est \emph{affine} d'anneau $A$. Comme $Z$ est alors réunion
finie d'ensembles de la forme $U\cap\complement V$, où $U$ et $V$ sont des
ouverts quasi-compacts (0\textsubscript{III}, 9.1.3), on voit, en considérant
encore une somme convenable de préschémas, qu'on peut se ramener au cas où
$Z=U\cap\complement V$~; comme d'ailleurs $U$ est réunion finie d'ouverts
affines, on peut même se borner au cas où $U$ est \emph{affine}. Comme $V$
est quasi-compact, il est réunion finie d'ouverts de la forme $X_{f_i}$, où
$f_i\in A$~; soit $\mathfrak J$ l'idéal de $A$ engendré par les $f_i$, et
soit $Z'$ le sous-schéma fermé affine de $X$ qu'il définit (et qui est par
construction de présentation finie sur $X$)~; on a par définition
$V=\complement Z'$, donc $U\cap\complement V=U\cap Z'$. Considérons le
schéma affine $X'=Z'\times_XU$

\oldpage[IV-1]{239}

et soit $f:X'\to X$ le morphisme structural~; on vient de voir que
l'immersion canonique $Z'\to X$ est de présentation finie, et il en est de
même de l'immersion ouverte $U\to X$, qui est quasi-compacte puisque $U$ et
$X$ sont affines~; on conclut donc de (1.6.2, (iv)) que $f$ est de
présentation finie, et l'on a $f(X')=Z'\cap U$ (I, 3.4.8), ce qui achève la
démonstration.
\end{proof}

\begin{theorem}[1.8.4]
\label{IV.1.8.4-fr}
\textup{(Chevalley).} --- Soit $f:X\to Y$ un morphisme de présentation
finie. Pour toute partie localement constructible $Z$ de $X$, $f(Z)$ est
localement constructible dans $Y$.
\end{theorem}

\begin{proof}
Soient $y\in Y$ et $V$ un voisinage ouvert affine de $y$~; comme $f$ est
quasi-compact et quasi-séparé, il en est de même de sa restriction
$f^{-1}(V)\to V$, donc $f^{-1}(V)$ est un préschéma quasi-compact et
quasi-séparé (1.2.3)~; comme $Z\cap f^{-1}(V)$ est constructible dans
$f^{-1}(V)$ (1.8.1), on voit qu'on peut se borner au cas où $Y$ est affine
et $Z$ constructible~; $X$ est alors quasi-compact et quasi-séparé, donc
(1.8.3), il existe un morphisme de présentation finie $g:X'\to X$ tel que
$Z=g(X')$~; comme $f\circ g$ est de présentation finie, on voit qu'on est
ramené au cas où $Z=X$~; autrement dit, il suffira de prouver le
\end{proof}

\begin{lemma}[1.8.4.1]
\label{IV.1.8.4.1-fr}
Soient $Y$ un schéma affine, $f:X\to Y$ un morphisme quasi-compact et
localement de présentation finie~; alors $f(X)$ est une partie constructible
de $Y$.
\end{lemma}

\begin{proof}
Comme $X$ est quasi-compact, il est réunion finie d'ouverts affines~; on peut
donc se borner au cas où $Y=\operatorname{Spec}(A)$,
$X=\operatorname{Spec}(B)$, $B$ étant une $A$-algèbre de présentation finie
(1.4.6). Or, on a le
\end{proof}

\begin{lemma}[1.8.4.2]
\label{IV.1.8.4.2-fr}
Soient $A_0$ un anneau, $(A_\lambda)_{\lambda\in L}$ un système inductif de
$A_0$-algèbres, $A=\varinjlim A_\lambda$~; si $B$ est une $A$-algèbre de
présentation finie, il existe un $\lambda$ et une $A_\lambda$-algèbre de
présentation finie $B_\lambda$ tels que $B$ soit isomorphe à
$B_\lambda\otimes_{A_\lambda}A$.
\end{lemma}

\begin{proof}
Par hypothèse, $B$ est isomorphe à une algèbre de la forme
$A[T_1,\ldots,T_n]/\mathfrak J$, où les $T_i$ sont des indéterminées et
$\mathfrak J$ un idéal de type fini~; soit $(F_j)$ ($1\leq j\leq m$) un
système de générateurs de $\mathfrak J$. Soit
$\varphi_\lambda:A_\lambda\to A$ l'homomorphisme canonique. Comme $L$ est
filtrant, il existe $\lambda$ tel que chacun des coefficients de chacun des
polynômes $F_j$ soit l'image par $\varphi_\lambda$ d'un élément de
$A_\lambda$. Autrement dit, il existe un système de polynômes
$(F_{j\lambda})_{1\leq j\leq m}$ de $A_\lambda[T_1,\ldots,T_n]$ tels que
$F_j=\varphi_\lambda(F_{j\lambda})$ pour $1\leq j\leq m$. Si
$\mathfrak J_\lambda$ est l'idéal de $A_\lambda[T_1,\ldots,T_n]$ engendré
par les $F_{j\lambda}$, $\mathfrak J$ est l'image de
$\mathfrak J_\lambda\otimes_{A_\lambda}A$ dans
$A[T_1,\ldots,T_n]=A_\lambda[T_1,\ldots,T_n]\otimes_{A_\lambda}A$~;
l'anneau $B_\lambda=A_\lambda[T_1,\ldots,T_n]/\mathfrak J_\lambda$ répond à
la question.
\end{proof}

On appliquera ici ce lemme en considérant $A$ comme limite inductive de ses
sous-$\mathbf Z$-algèbres \emph{de type fini}. On voit donc qu'il existe une
telle sous-$\mathbf Z$-algèbre $A_0$ de $A$, et une $A_0$-algèbre de
présentation finie $B_0$ telles que $B$ soit isomorphe à
$B_0\otimes_{A_0}A$~; posons $X_0=\operatorname{Spec}(B_0)$,
$Y_0=\operatorname{Spec}(A_0)$, de sorte que l'on a $X=X_0\times_{Y_0}Y$,
$f$ étant la projection $X\to Y$~; soient $f_0:X_0\to Y_0$,
$g_0:Y\to Y_0$ les morphismes structuraux~; il résulte de (I, 3.4.8) que
l'on a $f(X)=g_0^{-1}(f_0(X_0))$~; compte tenu de (1.8.2), il suffit donc de
montrer que $f_0(X_0)$ est constructible. Autrement dit, on est finalement
ramené à démontrer le

\begin{corollary}[1.8.5]
\label{IV.1.8.5-fr}
Soient $Y$ un préschéma noethérien, $f:X\to Y$ un morphisme de type fini.
Pour toute partie constructible $Z$ de $X$, $f(Z)$ est une partie
constructible de $Y$.
\end{corollary}

\oldpage[IV-1]{240}

\begin{proof}
On se ramène comme ci-dessus à prouver que $f(X)$ est constructible.
Appliquons le critère (0\textsubscript{III}, 9.2.3)~: il faut prouver que pour
toute partie fermée irréductible $T$ de $Y$, $T\cap f(X)=f(f^{-1}(T))$ est
rare dans $T$ ou contient un ouvert non vide de $T$~; en prenant un
sous-préschéma de $Y$ ayant $T$ pour espace sous-jacent (I, 5.2.1) et en
considérant son image réciproque par $f$, on voit finalement (compte tenu de
(1.5.4, (iii)) qu'on est ramené à prouver que si l'on suppose $Y$
\emph{irréductible et $f$ dominant}, $f(X)$ contient un ouvert non vide de
$Y$. On peut en outre supposer $Y$ affine~; alors $X$ est réunion finie
d'ouverts affines $V_i$, et comme $f(X)$ est dense dans $Y$, il en est de
même d'un au moins des $f(V_i)$. On peut donc aussi supposer $X$ affine~; si
$(X_j)$ est la famille (finie) des composantes irréductibles de $X$, on voit
comme ci-dessus qu'un au moins des $f(X_j)$ est dense dans $Y$~; on peut donc
supposer $X$ irréductible. Enfin, en remplaçant $f$ par $f_{\mathrm{red}}$
(1.5.4, (vi)) on peut supposer $X$ et $Y$ \emph{intègres}. Alors
$X=\operatorname{Spec}(B)$, $Y=\operatorname{Spec}(A)$, où $A$ est un
anneau intègre noethérien, $B$ une $A$-algèbre intègre de type fini contenant
$A$ (I, 1.2.7). Il suffit de montrer qu'il existe $g\in A$ tel que, pour tout
$y\in D(g)$ (c'est-à-dire tel que $g\notin\mathfrak p_y$) l'idéal
$\mathfrak p_y$ soit l'intersection de $A$ et d'un idéal premier de $B$, car
cela montrera que $D(g)\subset f(X)$. Finalement il suffit de prouver le
\end{proof}

\begin{lemma}[1.8.5.1]
\label{IV.1.8.5.1-fr}
Soient $A$ un anneau intègre, $B$ une $A$-algèbre intègre de type fini. Il
existe $g\in A$ tel que tout homomorphisme de $A$ dans un corps
algébriquement clos $\Omega$, non nul en $g$, se prolonge en un homomorphisme
de $B$ dans $\Omega$.
\end{lemma}

\begin{proof}
Or c'est là un résultat classique d'algèbre commutative (Bourbaki,
\emph{Alg. comm.}, chap. V, \S~3, n\textsuperscript{o}~1, cor. 3 du th. 1).
\end{proof}

\begin{corollary}[1.8.6]
\label{IV.1.8.6-fr}
Soient $Y$ un préschéma irréductible, $\eta$ son point générique,
$f:X\to Y$ un morphisme localement de type fini. Si $f^{-1}(\eta)$ n'est pas
vide, il existe un voisinage ouvert $U$ de $\eta$ dans $Y$ tel que
$f^{-1}(y)$ soit non vide pour tout $y\in U$. Si $\mathscr F$ est un
$\mathcal O_X$-Module quasi-cohérent de type fini, et si
$\mathscr F_\eta=\mathscr F\otimes_{\mathcal O_Y}k(\eta)$ n'est pas nul,
alors il existe un voisinage ouvert $U'$ de $\eta$ dans $Y$ tel que
$\mathscr F_y=\mathscr F\otimes_{\mathcal O_Y}k(y)$ ne soit pas nul pour
tout $y\in U'$.
\end{corollary}

\begin{proof}
Si $p_y$ est la projection canonique de la fibre
$f^{-1}(y)=X\times_Y\operatorname{Spec}(k(y))$ dans $X$, on a
$\mathscr F_y=p_y^*(\mathscr F)$, donc
$\operatorname{Supp}(\mathscr F_y)=p_y^{-1}(\operatorname{Supp}(\mathscr F))
=\operatorname{Supp}(\mathscr F)\cap f^{-1}(y)$ en vertu de (I, 9.1.13.1) et
(I, 3.6.1), puisque $\mathscr F$ est de type fini~; en outre
$\operatorname{Supp}(\mathscr F)$ est fermé dans $X$ (0\textsubscript{I},
5.2.2), et si $Z$ est un sous-préschéma fermé de $X$ ayant pour espace
sous-jacent $\operatorname{Supp}(\mathscr F)$ (I, 5.2.1), et $j$ l'immersion
canonique $Z\to X$, $f\circ j$ est localement de type fini (1.3.4)~; cela
prouve que la première assertion entraîne la seconde. Pour prouver la
première assertion, remarquons d'abord que l'on peut supposer $X$ et $Y$
réduits (1.3.4, (vi)), autrement dit $Y$ intègre. Soit $\xi$ un point de
$f^{-1}(\eta)$~; on peut remplacer $X$ et $Y$ par des voisinages ouverts
affines de $\xi$ et $\eta$ respectivement, et par suite supposer que
$X=\operatorname{Spec}(B)$, $Y=\operatorname{Spec}(A)$, $B$ étant une
$A$-algèbre de type fini. En outre, si $Z$ est le sous-préschéma fermé réduit
de $X$ ayant $\overline{\{\xi\}}$ pour ensemble sous-jacent (I, 5.2.1), on
peut, comme ci-dessus, remplacer $X$ par $Z$, autrement dit supposer $B$
\emph{intègre} (I, 5.1.4). Comme le morphisme $f$ est alors dominant,
l'homomorphisme correspondant $\varphi:A\to B$ est injectif (I, 1.2.7)~;
donc le corollaire résulte finalement du lemme (1.8.5.1).
\end{proof}

\oldpage[IV-1]{241}

\begin{proposition}[1.8.7]
\label{IV.1.8.7-fr}
Soient $S$ un préschéma, $g:X\to S$, $h:Y\to S$ deux morphismes de
présentation finie, $f:X\to Y$ un $S$-morphisme. Pour tout $s\in S$, posons
$X_s=g^{-1}(s)=X\times_S\operatorname{Spec}(k(s))$,
$Y_s=h^{-1}(s)=Y\times_S\operatorname{Spec}(k(s))$,
$f_s=f\times 1:X_s\to Y_s$. Alors l'ensemble des $s\in S$ tels que $f_s$
soit surjectif (resp. radiciel) est localement constructible.
\end{proposition}

\begin{proof}
Soit $E$ l'ensemble des $s\in S$ tels que $f_s$ soit surjectif~; l'ensemble
$S-E$ est égal à $h(Y-f(X))$~; or $f$ est de présentation finie (1.6.2, (v))
donc $f(X)$ est localement constructible dans $Y$ (1.8.4)~; comme $h$ est de
présentation finie, $h(Y-f(X))$ est localement constructible dans $S$, donc
aussi $E$.

Pour montrer que l'ensemble $E'$ des $s\in S$ tels que $f_s$ soit radiciel
est localement constructible, nous utiliserons le
\end{proof}

\begin{lemma}[1.8.7.1]
\label{IV.1.8.7.1-fr}
Soient $U$, $V$ deux préschémas~; pour qu'un morphisme $f:U\to V$ soit
radiciel, il faut et il suffit que le morphisme diagonal
$\Delta_f:U\to U\times_VU$ soit surjectif. Par suite, tout morphisme radiciel
est séparé.
\end{lemma}

\begin{proof}
En effet, $\Delta_f$ étant une immersion (I, 5.3.9), est un morphisme
localement de type fini (1.3.4)~; dire que $\Delta_f$ est surjectif signifie
donc que pour tout corps algébriquement clos $K$, l'application
correspondante $U(K)\to(U\times_VU)(K)=U(K)\times_{V(K)}U(K)$ est surjective
(1.3.7 et I, 3.4.2.1)~; mais par définition d'un produit fibré d'ensembles,
cela signifie que l'application $U(K)\to V(K)$ correspondant à $f$ est
injective, et cela équivaut à dire que $f$ est radiciel (I, 3.5.5).

Cela étant, $\Delta_{f_s}$ se déduit de $\Delta_f:X\to X\times_YX$ par le
changement de base $\operatorname{Spec}(k(s))\to S$ (I, 5.3.4). Comme $f$
est de présentation finie, il en est de même du morphisme structural
$X\times_YX\to Y$ (1.6.2, (iv)), donc aussi de $\Delta_f$ (1.6.2, (v))~; il
suffit donc d'appliquer la première partie de la proposition à $\Delta_f$,
en utilisant le lemme (1.8.7.1).
\end{proof}

\subsection{Ensembles pro-constructibles et ind-constructibles}
\label{subsection:IV.1.9-fr}

\begin{lemma}[1.9.1]
\label{IV.1.9.1-fr}
Soit $(A_\alpha)_{\alpha\in I}$ un système inductif d'anneaux dont l'ensemble
d'indices $I$ est filtrant à droite, et soit $A=\varinjlim A_\alpha$. Pour que
$A=0$, il faut et il suffit qu'il existe un indice $\gamma$ tel que
$A_\gamma=0$ (et l'on a alors $A_\beta=0$ pour tout $\beta\geq\gamma$).
\end{lemma}

\begin{proof}
En effet, pour tout $\alpha\in I$, l'homomorphisme canonique
$\varphi_\alpha:A_\alpha\to A$ transforme l'élément unité en élément unité~;
dire que $A=0$ signifie que $\varphi_\alpha(1)=0$, ou encore
$\varphi_\alpha(1)=\varphi_\alpha(0)$~; on sait que cela équivaut à dire
qu'il existe un indice $\gamma\geq\alpha$ tel que l'homomorphisme
$\varphi_{\gamma\alpha}:A_\alpha\to A_\gamma$ soit tel que
$\varphi_{\gamma\alpha}(1)=\varphi_{\gamma\alpha}(0)$, et cette dernière
relation équivaut à $A_\gamma=0$, d'où le lemme.
\end{proof}

\begin{proposition}[1.9.2]
\label{IV.1.9.2-fr}
Soient $B$ un anneau, $(A_\alpha)_{\alpha\in I}$ un système inductif de
$B$-algèbres dont l'ensemble d'indices $I$ est filtrant à droite, et soit
$A=\varinjlim A_\alpha$~; posons $Y=\operatorname{Spec}(B)$,
$X_\alpha=\operatorname{Spec}(A_\alpha)$, $X=\operatorname{Spec}(A)$, et
soient $u:X\to Y$, $u_\alpha:X_\alpha\to Y$ les morphismes correspondant aux
homomorphismes structuraux $B\to A$, $B\to A_\alpha$. Alors~:
\end{proposition}

\begin{env}[1.9.2.i]
\label{IV.1.9.2.i-fr}
\emph{Pour que $X=\varnothing$, il faut et il suffit qu'il existe un
$\gamma\in I$ tel que $X_\gamma=\varnothing$, et l'on a alors
$X_\alpha=\varnothing$ pour $\alpha\geq\gamma$.}
\end{env}

\oldpage[IV-1]{242}

\begin{env}[1.9.2.ii]
\label{IV.1.9.2.ii-fr}
\emph{On a}
\[
  u(X)=\bigcap_{\alpha\in I}u_\alpha(X_\alpha). \tag{1.9.2.1}
  \label{IV.1.9.2.1-fr}
\]
\end{env}

L'assertion (i) n'est autre que la traduction de (1.9.1). Pour démontrer
(ii), notons que, puisque $u$ se factorise en
$X\to X_\alpha\xrightarrow{u_\alpha}Y$, le premier membre de (1.9.2.1) est
contenu dans le second. Inversement, soit $y\in Y-u(X)$~; on a
$u^{-1}(y)=\varnothing$, et $u^{-1}(y)$ est l'espace sous-jacent au
$k(y)$-préschéma $\operatorname{Spec}(A\otimes_Bk(y))$~; comme dans la
catégorie des $B$-modules le foncteur $\varinjlim$ commute au produit
tensoriel, $A\otimes_Bk(y)$ est limite inductive du système inductif d'anneaux
$A_\alpha\otimes_Bk(y)$~; le lemme (1.9.1) montre qu'il existe $\alpha$ tel
que $A_\alpha\otimes_Bk(y)=0$, c'est-à-dire
$u_\alpha^{-1}(y)=\varnothing$, donc $y\notin u_\alpha(X_\alpha)$. C.Q.F.D.

\begin{proposition}[1.9.3]
\label{IV.1.9.3-fr}
Soient $X$ un préschéma quasi-compact et quasi-séparé, $E$ une partie de $X$,
$(X_\alpha)_{\alpha\in I}$ un recouvrement ouvert affine de $X$, et pour
tout $\alpha\in I$, posons $E_\alpha=E\cap X_\alpha$. Les conditions
suivantes sont équivalentes~:
\begin{enumerate}
  \item[a)] Il existe un morphisme quasi-compact $f:X'\to X$ tel que
  $E=f(X')$.
  \item[a$'$)] Pour tout $\alpha$, il existe un morphisme quasi-compact
  $f_\alpha:X'_\alpha\to X_\alpha$ tel que
  $E_\alpha=f_\alpha(X'_\alpha)$.
  \item[a$''$)] Pour tout $\alpha$, il existe un schéma affine $X''_\alpha$
  et un morphisme $f'_\alpha:X''_\alpha\to X_\alpha$ tel que
  $f'_\alpha(X''_\alpha)=E_\alpha$.
  \item[b)] $E$ est intersection de parties constructibles de $X$.
  \item[b$'$)] $F=X-E$ est réunion de parties constructibles de $X$.
\end{enumerate}
\end{proposition}

\begin{proof}
Il est clair que $b)$ et $b')$ sont équivalentes. La condition $a)$ entraîne
$a')$ en prenant pour $X'_\alpha$ le préschéma induit sur l'ouvert
$f^{-1}(X_\alpha)$ de $X'$ et pour $f_\alpha$ la restriction de $f$. Pour
voir que $a')$ entraîne $a'')$, il suffit de remarquer que $X'_\alpha$ est
réunion d'ouverts affines $X''_{\alpha\lambda}$ ($\lambda\in L$)~; soit
$X''_\alpha$ le préschéma \emph{somme} des $X''_{\alpha\lambda}$
($\lambda\in L$), qui est un schéma affine~; si
$g_\alpha:X''_\alpha\to X'_\alpha$ est le morphisme coïncidant avec
l'identité dans chacun des $X''_{\alpha\lambda}$, il est clair que si l'on
pose $f'_\alpha=f_\alpha\circ g_\alpha$, on a
$E_\alpha=f'_\alpha(X''_\alpha)$ puisque $g_\alpha$ est surjectif. Pour
montrer que $a'')$ entraîne $a)$, notons qu'il y a une partie finie $J$ de
$I$ telle que les $X_\alpha$ d'indices $\alpha\in J$ recouvrent $X$~; soit
$X'$ le préschéma somme des $X''_\alpha$ pour $\alpha\in J$, et soit
$f:X'\to X$ le morphisme coïncidant avec $j_\alpha\circ f'_\alpha$ pour tout
$\alpha\in J$, où $j_\alpha:X_\alpha\to X$ est l'injection canonique. Comme
$E$ est réunion des $E\cap X_\alpha$ pour $\alpha\in J$, on a bien
$E=f(X')$ et il reste à voir que $f$ est un morphisme quasi-compact~; mais
l'hypothèse que $X$ est quasi-séparé entraîne que $j_\alpha$ est
quasi-compact (1.2.7), donc il en est de même de
$j_\alpha\circ f'_\alpha$ (1.1.1 et 1.1.2) et finalement de $f$ (1.1.6).

Reste donc à prouver l'équivalence de $a'')$ et de $b)$. Prouvons en premier
lieu que $a'')$ entraîne $b)$~: considérons une partie finie $J$ de $I$ telle
que les $X_\alpha$ pour $\alpha\in J$ recouvrent $X$~; il suffira de montrer
que, pour tout $\alpha\in J$, $E_\alpha$ est intersection de parties
constructibles $F_{\alpha\lambda}$ de $X_\alpha$ ($\lambda\in L_\alpha$)~;
en effet, tout $F_{\alpha\lambda}$ est aussi une partie constructible de $X$
en vertu de (0\textsubscript{III}, 9.1.8, (ii)), car l'hypothèse que $X$ est
quasi-séparé entraîne que tout $X_\alpha$ est rétrocompact dans $X$ (1.2.7)~;
$E$ étant réunion des $E_\alpha$ pour $\alpha\in J$, est intersection des
réunions (finies) $\bigcup_{\alpha\in J}F_{\alpha,\lambda(\alpha)}$, pour
tous les choix des $\lambda(\alpha)\in L_\alpha$~; ces réunions étant
constructibles dans $X$, cela établit notre assertion. On peut donc se borner
au cas

\oldpage[IV-1]{243}

où $X=\operatorname{Spec}(A)$ et où $E=f(X')$, où
$X'=\operatorname{Spec}(A')$ est affine, $f:X'\to X$ un morphisme
correspondant à un homomorphisme d'anneaux $A\to A'$. Notons maintenant que
$A'$ est limite inductive de la famille $(A'_\lambda)$ de ses
sous-$A$-algèbres \emph{de type fini}, ordonnée par inclusion~; si
$X'_\lambda=\operatorname{Spec}(A'_\lambda)$, $f$ se factorise en
$X'\to X'_\lambda\xrightarrow{f_\lambda}X$ et il résulte de (1.9.2, (ii))
que l'on a $f(X')=\bigcap_\lambda f_\lambda(X'_\lambda)$~; si l'on démontre
que $f_\lambda(X'_\lambda)$ est intersection de parties constructibles de
$X$, il en sera de même de $f(X')$. On peut donc se borner au cas où $A'$ est
une $A$-algèbre \emph{de type fini}. Or, on a le lemme suivant~:
\end{proof}

\begin{lemma}[1.9.3.1]
\label{IV.1.9.3.1-fr}
Soient $A$ un anneau, $A'$ une $A$-algèbre de type fini. Alors $A'$ est
$A$-isomorphe à une limite inductive filtrante $\varinjlim A''_\mu$ où les
$A''_\mu$ sont des $A$-algèbres de présentation finie (1.4.1).
\end{lemma}

\begin{proof}
En effet, on peut écrire $A'=B/\mathfrak J$, avec
$B=A[T_1,\ldots,T_n]$ et $\mathfrak J$ un idéal de $B$. Mais $\mathfrak J$
est limite inductive des idéaux \emph{de type fini} $\mathfrak J'_\mu$ de
$B$, contenus dans $\mathfrak J$, ordonnés par inclusion~; comme dans la
catégorie des $B$-modules le foncteur $\varinjlim$ est exact, on en conclut
que $A'=\varinjlim B/\mathfrak J'_\mu$ à un $A$-isomorphisme près.
\end{proof}

Le même raisonnement que ci-dessus permet alors de se ramener au cas où $A'$
est une $A$-algèbre de présentation finie~; mais alors $f(X')$ est
constructible dans $X$ en vertu du théorème de Chevalley (1.8.5), ce qui
prouve $b)$.

Démontrons enfin que $b)$ entraîne $a'')$. Si $E$ est intersection d'une
famille $(G_\lambda)_{\lambda\in L}$ de parties constructibles de $X$,
chaque $E_\alpha$ est intersection des $G_\lambda\cap X_\alpha$, qui sont
constructibles dans $X_\alpha$ (0\textsubscript{III}, 9.1.8, (i)), donc on
est ramené au cas où $X=\operatorname{Spec}(A)$ est affine.

On sait alors (1.8.3) que pour tout $\lambda\in L$ il existe un morphisme
$f_\lambda:X'_\lambda\to X$, où
$X'_\lambda=\operatorname{Spec}(A'_\lambda)$ est affine, tel que
$f_\lambda(X'_\lambda)=G_\lambda$. Il suffit donc de prouver le lemme
suivant~:

\begin{lemma}[1.9.3.2]
\label{IV.1.9.3.2-fr}
Soient $(A'_\lambda)_{\lambda\in L}$ une famille de $A$-algèbres,
$X=\operatorname{Spec}(A)$, $X'_\lambda=\operatorname{Spec}(A'_\lambda)$,
et soit $f_\lambda:X'_\lambda\to X$ le morphisme structural. Pour toute
partie finie $J$ de $L$, posons
$A'_J=\bigotimes_{\lambda\in J}A'_\lambda$ (produit tensoriel de
$A$-algèbres), $X'_J=\operatorname{Spec}(A'_J)$, et soit $f_J:X'_J\to X$ le
morphisme structural. On a alors
$f_J(X'_J)=\bigcap_{\lambda\in J}f_\lambda(X'_\lambda)$. Si
$A'=\varinjlim A'_J$, $J$ parcourant l'ensemble filtrant des parties finies
de $L$, $X'=\operatorname{Spec}(A')$, et si $f:X'\to X$ est le morphisme
structural, on a $f(X')=\bigcap_{\lambda\in L}f_\lambda(X'_\lambda)$.
\end{lemma}

\begin{proof}
La première assertion résulte de (I, 3.4.7)~; la deuxième résulte de (1.9.2,
(ii)), qui donne la relation $f(X')=\bigcap_Jf_J(X'_J)$.
\end{proof}

\begin{definition}[1.9.4]
\label{IV.1.9.4-fr}
Soit $X$ un espace topologique. On dit qu'une partie $E$ de $X$ est
\emph{pro-constructible} (resp. \emph{ind-constructible}) dans $X$ si, pour
tout $x\in X$, il existe un voisinage ouvert $U$ de $x$ dans $X$ tel que
$E\cap U$ soit intersection (resp. réunion) de parties localement
constructibles de $U$.
\end{definition}

Tenant compte de (1.2.7) et de (0\textsubscript{III}, 9.1.8 et 9.1.12), les
conditions équivalentes de la proposition (1.9.3) s'expriment donc en disant
que $E$ est \emph{pro-constructible} dans $X$, les ensembles localement
constructibles dans $X$ étant identiques aux ensembles constructibles dans
$X$ sous les hypothèses de (1.9.3).

\oldpage[IV-1]{244}

\begin{proposition}[1.9.5]
\label{IV.1.9.5-fr}
Soit $X$ un préschéma.
\begin{enumerate}
  \item[(i)] Pour qu'une partie $E$ de $X$ soit pro-constructible, il faut et
  il suffit que $X-E$ soit ind-constructible.
  \item[(ii)] Toute réunion finie et toute intersection de parties
  pro-constructibles de $X$ sont pro-constructibles. Toute intersection
  finie et toute réunion de parties ind-constructibles de $X$ sont
  ind-constructibles.
  \item[(iii)] Toute intersection (resp. réunion) de parties localement
  constructibles de $X$ est pro-constructible (resp. ind-constructible).
  Inversement, si $X$ est quasi-compact et quasi-séparé, toute partie
  pro-constructible (resp. ind-constructible) de $X$ est intersection (resp.
  réunion) de parties constructibles de $X$.
  \item[(iv)] Soit $(U_\alpha)$ un recouvrement ouvert de $X$. Pour qu'une
  partie $E$ de $X$ soit pro-constructible (resp. ind-constructible) dans
  $X$, il faut et il suffit que pour tout $\alpha$, $E\cap U_\alpha$ soit
  pro-constructible (rep. ind-constructible) dans $U_\alpha$.
  \item[(v)] Toute partie pro-constructible $E$ de $X$ est rétrocompacte
  dans $X$. Pour qu'une partie localement fermée $E$ de $X$ soit
  rétrocompacte dans $X$, il faut et il suffit qu'elle soit
  pro-constructible dans $X$.
  \item[(vi)] Soit $f:X'\to X$ un morphisme de préschémas~; pour toute partie
  pro-constructible (resp. ind-constructible) $E$ de $X$, $f^{-1}(E)$ est
  pro-constructible (resp. ind-constructible) dans $X'$.
  \item[(vii)] Soit $f:X'\to X$ un morphisme quasi-compact~; pour toute
  partie pro-constructible $E'$ de $X'$, $f(E')$ est pro-constructible dans
  $X$~; en particulier $f(X')$ est pro-constructible dans $X$.
  \item[(viii)] Soit $f:X'\to X$ un morphisme localement de présentation
  finie~; pour toute partie ind-constructible $E'$ de $X'$, $f(E')$ est
  ind-constructible dans $X$~; en particulier $f(X')$ est ind-constructible
  dans $X$.
  \item[(ix)] Supposons que $X$ soit quasi-compact et quasi-séparé~; alors,
  pour qu'une partie $E$ de $X$ soit pro-constructible, il faut et il suffit
  qu'il existe un schéma affine $X'$ et un morphisme (nécessairement
  quasi-compact) $f:X'\to X$ tels que $E=f(X')$.
\end{enumerate}
\end{proposition}

\begin{proof}
Les propriétés (i), (ii), (iv) et la première assertion de (iii) résultent
de la définition (1.9.4) et sont valables pour un espace topologique
quelconque, en utilisant la distributivité mutuelle de l'intersection et de
la réunion et le fait que si $T$ est localement constructible dans $X$,
$T\cap U$ est localement constructible dans $U$ pour tout ouvert $U$ de $X$.
Si $X$ est quasi-compact et quasi-séparé et $E$ est pro-constructible dans
$X$, il y a un recouvrement fini $(U_i)$ de $X$ formé d'ouverts affines tels
que $E\cap U_i$ soit intersection de parties constructibles $M_{i\lambda}$
de $U_i$ ($\lambda\in L_i$)~; les $M_{i\lambda}$ sont aussi constructibles
dans $X$ en vertu de (1.2.7) et (0\textsubscript{III}, 9.1.8, (ii)) et $E$ est
intersection des réunions finies $\bigcup_iM_{i,\lambda(i)}$ (où
$\lambda(i)\in L_i$ pour tout $i$), qui sont des parties constructibles dans
$X$~; cela démontre la seconde assertion de (iii). L'assertion (ix) résulte
de (iii) et de (1.9.3). Pour démontrer la première assertion de (v), on peut
se limiter au cas où $X$ est affine, et prouver alors que $E$ est
quasi-compact (0\textsubscript{III}, 9.1.1)~; mais comme $X$ est alors
quasi-séparé, il existe en vertu de (ix) un morphisme quasi-compact
$f:X'\to X$ tel que $E=f(X')$~; comme $X'$ est quasi-compact, il en est de
même de son image $E$ par une application continue.

Pour prouver (vi), on peut se borner, en vertu de (iv), au cas où $X$ est
affine~; la conclusion résulte alors de (iii) et de (1.8.2).

\oldpage[IV-1]{245}

De même, pour démontrer (vii), on peut se borner au cas où $X$ est affine~;
alors $X'$ est réunion finie d'ouverts affines $X'_i$ et $f(E')$ est réunion
des $f(E'\cap X'_i)$, si bien qu'on peut aussi supposer que $X'$ est affine,
en vertu de (ii)~; mais alors $E'=g(X'')$, où $g:X''\to X'$ est un morphisme
quasi-compact, en vertu de (ix), et l'on a $f(E')=f(g(X''))$, donc $f(E')$ est
pro-constructible puisque $g\circ f$ est quasi-compact.

On déduit de (vii) la démonstration de la seconde assertion de (v)~: en
effet, soit $E$ une partie localement fermée et rétrocompacte dans $X$, et
considérons un sous-préschéma de $X$ ayant $E$ pour espace sous-jacent (I,
5.2.1)~; l'injection canonique $j:E\to X$ est alors par hypothèse un
morphisme quasi-compact, et il résulte de (vii) que $E=j(E)$ est
pro-constructible dans $X$.

Enfin, pour prouver (viii), on peut encore supposer $X$ affine~; en outre, si
$(X'_\alpha)$ est un recouvrement ouvert affine de $X'$ tel que
$E'\cap X'_\alpha$ soit réunion de parties constructibles de $X'_\alpha$
((iii) et (iv)), $f(E')$ est réunion des $f(E'\cap X'_\alpha)$ et, en vertu
du théorème de Chevalley (1.8.4), chacune des $f(E'\cap X'_\alpha)$ est
réunion de parties constructibles de $X$, d'où la conclusion.
\end{proof}

\begin{env}[1.9.6]
\label{IV.1.9.6-fr}
\textbf{Exemples} (1.9.6). --- Toute partie \emph{finie} d'un préschéma $X$
est pro-constructible~: en effet, il suffit de considérer les parties
$\{x\}$ réduites à un seul point (1.9.5, (ii)) et $\{x\}$ est l'image de
$\operatorname{Spec}(k(x))$ par le morphisme canonique
$\operatorname{Spec}(k(x))\to X$, qui est quasi-compact~; d'où la conclusion
par (1.9.5, (vii)). Par contre, une partie $\{x\}$ réduite à un point n'est
pas nécessairement ind-constructible~; par exemple, soit $A$ un anneau
noethérien intègre ayant une infinité d'idéaux maximaux, et soit $x$ le point
générique de $X=\operatorname{Spec}(A)$~; si $\{x\}$ était ind-constructible
dans $X$, il serait constructible en vertu de (1.9.5, (iii)) et par suite
contiendrait un ouvert non vide de $X$ (0\textsubscript{III}, 9.2.2)~; mais
cela contredit l'hypothèse que $A$ possède une infinité d'idéaux maximaux, en
vertu du théorème d'Artin-Tate (0, 16.3.3).

Toute partie \emph{fermée} $Y$ d'un préschéma $X$ est pro-constructible, par
(1.9.5, (v)). Toute partie ouverte de $X$ est donc ind-constructible, mais
une partie ouverte $U$ de $X$ n'est pro-constructible que si elle est
\emph{rétrocompacte}, en vertu encore de (1.9.5, (v)).

Enfin, si $X$ est un préschéma \emph{noethérien}, donc quasi-séparé (1.2.8),
il résulte de (1.9.5, (iii)) et de (0\textsubscript{III}, 9.1.7) que les
parties ind-constructibles de $X$ sont les \emph{réunions de parties
localement fermées de $X$}.
\end{env}

\begin{theorem}[1.9.7]
\label{IV.1.9.7-fr}
Soient $X$ un préschéma quasi-compact, $E$ une partie ind-constructible de
$X$, $(F_\lambda)_{\lambda\in L}$ une famille de parties pro-constructibles
de $X$ telle que $\bigcap_{\lambda\in L}F_\lambda\subset E$~; alors il existe
une partie finie $J$ de $L$ telle que
$\bigcap_{\lambda\in J}F_\lambda\subset E$.
\end{theorem}

\begin{proof}
Notons que les ensembles $F=X-E$ et $F\cap F_\lambda$ sont
pro-constructibles, donc on est ramené au
\end{proof}

\begin{corollary}[1.9.8]
\label{IV.1.9.8-fr}
Soient $X$ un préschéma quasi-compact, $(F_\lambda)_{\lambda\in L}$ une
famille de parties pro-constructibles de $X$ dont l'intersection est vide~;
alors il existe une sous-famille finie $(F_\lambda)_{\lambda\in J}$ dont
l'intersection est vide.
\end{corollary}

\begin{proof}
Recouvrant $X$ par un nombre fini d'ouverts affines, et utilisant (1.9.5,
(iv)), on est ramené au cas où $X$ est affine~; en vertu de (1.9.5, (ix)), on
a $F_\lambda=f_\lambda(X'_\lambda)$,

\oldpage[IV-1]{246}

où $X'_\lambda=\operatorname{Spec}(A'_\lambda)$ est affine, et $f_\lambda$
est un morphisme $X'_\lambda\to X$~; alors, avec les notations de (1.9.3.2),
on a par hypothèse $f(X')=\varnothing$, donc $X'=\varnothing$~; mais cela
entraîne par (1.9.2, (i)) qu'il existe une partie finie $J$ de $L$ telle que
$X'_J=\varnothing$, donc
$f_J(X'_J)=\bigcap_{\lambda\in J}F_\lambda=\varnothing$.
\end{proof}

\begin{corollary}[1.9.9]
\label{IV.1.9.9-fr}
Soient $X$ un préschéma quasi-compact, $F$ une partie pro-constructible de
$X$, $(E_\lambda)_{\lambda\in L}$ une famille de parties ind-constructibles
de $X$ telle que $F\subset\bigcup_{\lambda\in L}E_\lambda$~; alors il existe
une partie finie $J$ de $L$ telle que
$F\subset\bigcup_{\lambda\in J}E_\lambda$. En particulier, de tout
recouvrement de $X$ par des parties ind-constructibles, on peut extraire un
recouvrement fini de $X$.
\end{corollary}

\begin{proof}
Cela résulte de (1.9.7) par passage aux complémentaires.
\end{proof}

\begin{proposition}[1.9.10]
\label{IV.1.9.10-fr}
Soit $X$ un préschéma. Pour qu'une partie $E$ de $X$ soit
ind-constructible, il est nécessaire que, pour tout $x\in X$, l'intersection
$E\cap\overline{\{x\}}$ soit un voisinage de $x$ dans
$\overline{\{x\}}$. Si $X$ est localement noethérien, cette condition est
suffisante.
\end{proposition}

\begin{proof}
Supposons $E$ ind-constructible, et soit $Y$ l'intersection de
$\overline{\{x\}}$ et d'un ouvert affine dans $X$, contenant $x$~; il y a
donc un sous-préschéma de $X$ ayant $Y$ pour espace sous-jacent (I, 5.2.1),
et si $j:Y\to X$ est l'injection canonique,
$E\cap Y=j^{-1}(E)$ est ind-constructible dans $Y$ (1.9.5, (vi)). Comme $Y$
est quasi-compact et séparé, $E\cap Y$ est donc réunion de parties
constructibles de $Y$ (1.9.5, (iii))~; par suite, il y a deux ouverts $U$,
$V$ dans $Y$ tels que
$x\in U\cap\complement V\subset E\cap Y$ (0\textsubscript{III}, 9.1.3).
Mais comme $x$ est point générique de l'espace irréductible $Y$, $V$ est
nécessairement vide et l'on a $U\subset E\cap Y$.

Supposons maintenant $X$ localement noethérien, et soit $E$ une partie de
$X$ vérifiant la condition de l'énoncé. En vertu de la définition (1.9.4),
on peut se borner au cas où $X$ est noethérien. Alors, pour tout $x\in E$,
$E\cap\overline{\{x\}}$ contient une partie non vide de la forme
$U\cap\overline{\{x\}}$, où $U$ est ouvert dans $X$~; cela montre que $E$
est réunion de parties localement fermées de $X$, donc est
ind-constructible (1.9.6).
\end{proof}

\begin{proposition}[1.9.11]
\label{IV.1.9.11-fr}
Soient $X$ un préschéma, $E$ une partie de $X$. Les propriétés suivantes sont
équivalentes~:
\begin{enumerate}
  \item[a)] $E$ est localement constructible.
  \item[b)] $E$ est ind-constructible et pro-constructible.
  \item[c)] $E$ et $X-E$ sont pro-constructibles.
  \item[d)] $E$ et $X-E$ sont ind-constructibles.
\end{enumerate}
\end{proposition}

\begin{proof}
Il suffit évidemment de prouver que $d)$ entraîne $a)$, et on peut se borner
au cas où $X$ est affine. Alors (1.9.5, (iii)), $E$ (resp. $X-E$) est réunion
d'une famille $(E_\lambda)$ (resp. $(E'_\mu)$) de parties constructibles de
$X$~; comme les $E_\lambda$ et les $E'_\mu$ forment un recouvrement de $X$,
il résulte de (1.9.9) qu'il y a des indices en nombre fini $\lambda_i$,
$\mu_j$ tels que les $E_{\lambda_i}$ et les $E'_{\mu_j}$ forment un
recouvrement de $X$~; cela implique que $E$ est réunion des
$E_{\lambda_i}$, donc est constructible.
\end{proof}

\begin{corollary}[1.9.12]
\label{IV.1.9.12-fr}
Soit $f:X\to Y$ un morphisme surjectif, qui est, soit quasi-compact, soit
localement de présentation finie. Pour qu'une partie $E$ de $Y$ soit
localement constructible (resp. pro-constructible, resp. ind-constructible),
il faut et il suffit que $f^{-1}(E)$ le soit dans $X$.
\end{corollary}

\oldpage[IV-1]{247}

\begin{proof}
On sait que la condition est nécessaire ((1.8.2) et (1.9.5, (vi)))~; pour
montrer qu'elle est suffisante, on est ramené au cas où $X$ est affine, en
vertu de (1.9.5, (iv))~; en outre, en vertu de (1.9.11), on peut se borner au
cas où $f^{-1}(E)$ est pro-constructible, \emph{ou} au cas où $f^{-1}(E)$ est
ind-constructible. Si $f$ est surjectif et quasi-compact, et $f^{-1}(E)$
pro-constructible, il résulte de (1.9.5, (vii)) que
$E=f(f^{-1}(E))$ est pro-constructible. Si $f$ est surjectif et localement de
présentation finie, et $f^{-1}(E)$ ind-constructible,
$E=f(f^{-1}(E))$ est ind-constructible par (1.9.5, (viii)), ce qui achève la
démonstration.
\end{proof}

\begin{env}[1.9.13]
\label{IV.1.9.13-fr}
Soient $X$ un préschéma, $\mathscr T$ sa topologie~; il résulte de (1.9.5,
(i) et (ii)) que les parties ind-constructibles (resp. pro-constructibles) de
$X$ sont les parties ouvertes (resp. fermées) pour une topologie sur $X$,
appelée \emph{topologie constructible} et que nous noterons
$\mathscr T^{\mathrm{cons}}$~; nous désignerons par $X^{\mathrm{cons}}$
l'ensemble $X$ muni de la topologie $\mathscr T^{\mathrm{cons}}$.
\end{env}

\begin{proposition}[1.9.14]
\label{IV.1.9.14-fr}
Soient $X$ un préschéma, $\mathscr T$ sa topologie,
$\mathscr T^{\mathrm{cons}}$ la topologie constructible sur $X$.
\begin{enumerate}
  \item[(i)] La topologie $\mathscr T^{\mathrm{cons}}$ est plus fine que
  $\mathscr T$.
  \item[(ii)] Les parties localement constructibles de $X$ sont identiques
  aux parties à la fois ouvertes et fermées de l'espace $X^{\mathrm{cons}}$.
  \item[(iii)] Pour tout morphisme $f:X\to Y$, l'application sous-jacente de
  $X^{\mathrm{cons}}$ dans $Y^{\mathrm{cons}}$ est continue~; on la note
  $f^{\mathrm{cons}}$.
  \item[(iv)] Si le morphisme $f:X\to Y$ est quasi-compact,
  $f^{\mathrm{cons}}$ est une application fermée~; en particulier, si $f$
  est quasi-compact et bijectif, $f^{\mathrm{cons}}$ est un homéomorphisme.
  \item[(v)] Si un morphisme $f:X\to Y$ est localement de présentation
  finie, $f^{\mathrm{cons}}$ est une application ouverte~; en particulier,
  si $f$ est bijectif et localement de présentation finie,
  $f^{\mathrm{cons}}$ est un homéomorphisme.
  \item[(vi)] Pour tout ouvert $U$ de $X$, la topologie induite par
  $\mathscr T^{\mathrm{cons}}$ sur $U$ est identique à la topologie de
  $U^{\mathrm{cons}}$.
\end{enumerate}
\end{proposition}

\begin{proof}
En effet, (i) résulte de ce que tout ouvert pour $\mathscr T$ est un ouvert
pour $\mathscr T^{\mathrm{cons}}$ (1.9.6), et (ii) est une traduction de
(1.9.11)~; (iii), (iv) et (v) traduisent respectivement (1.9.5), (vi), (vii)
et (viii). Enfin pour démontrer (vi), il suffit de remarquer que l'injection
canonique $j:U\to X$ est localement de présentation finie (1.4.3), donc
$j^{\mathrm{cons}}:U^{\mathrm{cons}}\to X^{\mathrm{cons}}$ est une
application continue ouverte.
\end{proof}

\begin{proposition}[1.9.15]
\label{IV.1.9.15-fr}
Soit $X$ un préschéma.
\begin{enumerate}
  \item[(i)] Pour que $X^{\mathrm{cons}}$ soit quasi-compact, il faut et il
  suffit que $X$ soit quasi-compact.
  \item[(ii)] Pour que $X^{\mathrm{cons}}$ soit séparé, il faut et il suffit
  que $X$ soit quasi-séparé, et alors $X^{\mathrm{cons}}$ est localement
  compact et totalement discontinu.
  \item[(iii)] Pour que $X^{\mathrm{cons}}$ soit compact, il faut et il suffit
  que $X$ soit quasi-compact et quasi-séparé.
  \item[(iv)] Tout point de $X$ admet, pour la topologie
  $\mathscr T^{\mathrm{cons}}$, un voisinage ouvert et compact.
  \item[(v)] Pour qu'un morphisme $f:X\to Y$ soit quasi-compact, il faut et il
  suffit que l'application continue
  $f^{\mathrm{cons}}:X^{\mathrm{cons}}\to Y^{\mathrm{cons}}$ soit propre.
\end{enumerate}
\end{proposition}

\begin{proof}
(i) Comme $\mathscr T^{\mathrm{cons}}$ est plus fine que $\mathscr T$, il est
clair que si $X^{\mathrm{cons}}$ est quasi-compact, il en est de même de $X$~;
la réciproque résulte de (1.9.9).

(ii) Supposons $X$ quasi-séparé, et montrons que $X^{\mathrm{cons}}$ est
totalement discontinu~: en effet, si $x$, $y$ sont deux points distincts de
$X$ et si par exemple $x\notin\overline{\{y\}}$, il existe un
\oldpage[IV-1]{248}
voisinage ouvert affine $U$ de $X$ ne contenant pas $y$~; comme $X$ est
quasi-séparé, $U$ est rétrocompact dans $X$ (1.2.7), donc $U$ et $X-U$ sont
localement constructibles dans $X$, et par suite ouverts dans
$X^{\mathrm{cons}}$ en vertu de (1.9.11), d'où notre assertion. Comme sur tout
ouvert affine $U$ de $X$, la topologie induite par $\mathscr T^{\mathrm{cons}}$
est celle de $U^{\mathrm{cons}}$ en vertu de (1.9.14, (vi)), il résulte de ce
qui précède que $X^{\mathrm{cons}}$ est localement compact, puisque $U$ est
ouvert dans $X^{\mathrm{cons}}$ et que $U^{\mathrm{cons}}$ est compact. Reste à
prouver que si $X^{\mathrm{cons}}$ est séparé, alors $X$ est quasi-séparé~;
considérons en effet un ouvert affine $U$ de $X$~; la topologie induite sur
$U$ par $\mathscr T^{\mathrm{cons}}$ est celle de $U^{\mathrm{cons}}$
(1.9.14, (vi)), donc $U$ est compact pour cette topologie induite, puisque
$U^{\mathrm{cons}}$ est compact par la première partie du raisonnement. Si
$V$ est un second ouvert affine dans $X$, $U\cap V$ est donc une partie
compacte de l'espace séparé $X^{\mathrm{cons}}$, étant intersection de deux
parties compactes de cet espace~; comme la topologie induite par
$\mathscr T^{\mathrm{cons}}$ sur $U\cap V$ est celle de
$(U\cap V)^{\mathrm{cons}}$ (1.9.14, (vi)) et que l'application identique
$(U\cap V)^{\mathrm{cons}}\to U\cap V$ est continue, on en conclut que
$U\cap V$ est quasi-compact pour la topologie induite par $\mathscr T$~; $X$
est par suite quasi-séparé en vertu de (1.2.7).

(iii) est un corollaire immédiat de (i) et (ii).

(iv) Pour tout $x\in X$, un voisinage ouvert affine $U$ de $x$ pour
$\mathscr T$ est aussi un voisinage de $x$ pour $\mathscr T^{\mathrm{cons}}$
et il est compact en vertu de (iii) et de (1.9.14, (vi)), la topologie induite
sur $U$ par $\mathscr T^{\mathrm{cons}}$ étant identique à celle de
$U^{\mathrm{cons}}$.

(v) Supposons $f$ quasi-compact~; alors on sait déjà (1.9.14, (iv)) que
$f^{\mathrm{cons}}$ est une application fermée. Soit d'autre part $y$ un
point de $Y$, et posons
\[
  Z=f^{-1}(y)=X\times_Y\operatorname{Spec}(k(y));
\]
$Z$ est quasi-compact (1.1.2, (iii)) et comme le morphisme canonique
$p:Z\to X$ est injectif, il résulte de (i) et de ce que l'application
$p^{\mathrm{cons}}$ est continue, que la topologie induite sur $f^{-1}(y)$
par celle de $X^{\mathrm{cons}}$ fait de $f^{-1}(y)$ un espace quasi-compact.
Cela prouve que $f^{\mathrm{cons}}$ est une application propre (Bourbaki,
\emph{Top. gén.}, chap. I\textsuperscript{er}, 3\textsuperscript{e} éd.,
\S~10, n\textsuperscript{o}~2, th.~1). Inversement, supposons que
l'application continue $f^{\mathrm{cons}}$ soit propre, et soit $V$ un ouvert
quasi-compact de $Y$~; si $h:V\to Y$ est l'injection canonique,
$h^{\mathrm{cons}}:V^{\mathrm{cons}}\to Y^{\mathrm{cons}}$ est continue et
injective et $V^{\mathrm{cons}}$ est quasi-compact par (i), donc la topologie
induite sur $V$ par celle de $Y^{\mathrm{cons}}$ fait de $V$ un espace
quasi-compact. L'hypothèse que $f^{\mathrm{cons}}$ est propre entraîne alors
que la topologie induite sur $f^{-1}(V)$ par celle de $X^{\mathrm{cons}}$ fait
de $f^{-1}(V)$ un espace quasi-compact (\emph{loc. cit.}, prop.~6), donc
$f^{-1}(V)$ est aussi un sous-espace quasi-compact de $X$, ce qui montre que
le morphisme $f$ est quasi-compact.
\end{proof}

\begin{env}[1.9.16]
\label{IV.1.9.16-fr}
Nous montrerons plus tard comment on peut, pour tout préschéma $X$, munir
l'espace $X^{\mathrm{cons}}$ d'un faisceau d'anneaux qui en fait un préschéma
dont les anneaux locaux sont tous des corps, identiques aux corps résiduels aux
points de $X$. De tels préschémas s'introduisent par exemple de façon naturelle
dans la traduction, en langage des schémas, des constructions de
Néron [32] dans sa théorie de la réduction des variétés abéliennes.
\end{env}

\oldpage[IV-1]{249}

\subsection{Application aux morphismes ouverts}
\label{subsection:IV.1.10-fr}

\begin{theorem}[1.10.1]
\label{IV.1.10.1-fr}
Soient $X$ un préschéma, $E$ une partie ind-constructible de $X$ (1.9.4), $x$
un point de $X$. Pour que $x$ soit intérieur à $E$, il faut et il suffit que
toute générisation ($0_{\mathrm I}$, 2.1.2) $x'$ de $x$ appartienne à $E$,
autrement dit (I, 2.4.2) que l'on ait
$\operatorname{Spec}(\mathcal O_{X,x})\subset E$.
\end{theorem}

\begin{proof}
La condition est évidemment nécessaire, tout voisinage de $x$ contenant les
générisations de $x$. Pour voir qu'elle est suffisante, on peut évidemment se
borner au cas où $X=\operatorname{Spec}(A)$ est affine, $x$ étant un idéal premier
$\mathfrak p$ de $A$. On sait alors (1.9.5, (ix)) qu'il existe un schéma affine
$X'=\operatorname{Spec}(A')$ et un morphisme $f:X'\to X$ tels que
$f(X')=X-E$. Posons $Y=\operatorname{Spec}(\mathcal O_{X,x})$ et
$Y'=X'\times_XY$~; l'hypothèse signifie (en
vertu de (I, 3.6.5)) que l'on a $Y'=\varnothing$, autrement dit
$A'\otimes_AA_{\mathfrak p}=0$. Comme
$A_{\mathfrak p}=\varinjlim A_t$, où $t$ parcourt $A-\mathfrak p$
($0_{\mathrm I}$, 1.4.5), on a $\varinjlim A_t'=0$, le foncteur
$\varinjlim$ commutant au produit tensoriel. Par suite (1.9.1), il existe un
$t\in A-\mathfrak p$ tel que $A_t'=0$, et $\mathrm D(t)$ est alors un
voisinage ouvert de $x$ dans $X$ contenu dans $E$.
\end{proof}

\begin{corollary}[1.10.2]
\label{IV.1.10.2-fr}
Soient $f:X\to Y$ un morphisme localement de présentation finie, et soit
$y\in Y$. Pour que $y$ soit intérieur à $f(X)$ il faut et il suffit que toute
générisation $y'$ de $y$ appartienne à $f(X)$.
\end{corollary}

\begin{proof}
On sait en effet (1.9.5, (viii)) que $f(X)$ est ind-constructible dans $Y$ et
il suffit d'appliquer (1.10.1).
\end{proof}

Nous dirons qu'une application continue $f:X\to Y$ est \emph{ouverte en un
point} $x\in X$ si l'image par $f$ de \emph{tout} voisinage de $x$ dans $X$
est un voisinage de $f(x)$ dans $Y$.

\begin{proposition}[1.10.3]
\label{IV.1.10.3-fr}
Soient $f:X\to Y$ un morphisme, $x$ un point de $X$, $y=f(x)$. Considérons
les conditions suivantes~:
\begin{enumerate}
  \item[a)] $f$ est ouvert au point $x$.
  \item[b)] Pour toute générisation $y'$ de $y$, il existe une générisation
  $x'$ de $x$ telle que $f(x')=y'$.
  \item[b')] L'image par $f$ de $\operatorname{Spec}(\mathcal O_{X,x})$ est
  $\operatorname{Spec}(\mathcal O_{Y,y})$.
  \item[c)] Pour toute partie fermée irréductible $Y'$ de $Y$ contenant $y$,
  il existe une composante irréductible de $X'=f^{-1}(Y')$ contenant $x$ et
  dominant $Y'$.
\end{enumerate}
Alors on a les implications
$\mathrm{a)}\Rightarrow\mathrm{b)}\Leftrightarrow\mathrm{b')}
\Leftrightarrow\mathrm{c)}$. Si en outre $f$ est localement de présentation
finie, les quatre conditions sont équivalentes.
\end{proposition}

\begin{proof}
Par définition d'une générisation, l'image par $f$ de
$\operatorname{Spec}(\mathcal O_{X,x})$ est toujours contenue dans
$\operatorname{Spec}(\mathcal O_{Y,y})$, donc b) et b') sont équivalentes. Il
est immédiat
que c) entraîne b), car si $y'$ est une générisation de $y$,
$Y'=\overline{\{y'\}}$ est une partie irréductible de $Y$ contenant $y$~; si
$x'$ est le point générique d'une composante irréductible de $f^{-1}(Y')$ qui
contient $x$ et domine $Y'$, on a $f(x')=y'$ ($0_{\mathrm I}$, 2.1.5) et $x'$
est une générisation de $x$. Inversement, b) entraîne c)~: en effet, soit
$y'$ le point générique de $Y'$ et soit $x'$ une générisation de $x$ telle
que $f(x')=y'$~; soit $Z'$ une composante irréductible de $f^{-1}(Y')$
contenant $x'$ (donc $x$)~; comme $f(x')$ est déjà dense dans $Y'$, il en est
ainsi \emph{a fortiori} de $f(Z')$.

Pour voir que a) entraîne b'), on peut évidemment se borner au cas où
$X=\operatorname{Spec}(B)$ et $Y=\operatorname{Spec}(A)$ sont affines, $x$
étant un idéal premier
$\mathfrak q$ de $B$, de sorte que
$\mathcal O_{X,x}=B_{\mathfrak q}=\varinjlim B_t$,
\oldpage[IV-1]{250}
où $t$ parcourt $B-\mathfrak q$. On a, par suite (1.9.2),
\[
  f(\operatorname{Spec}(\mathcal O_{X,x}))
  =\bigcap_t f(\operatorname{Spec}(B_t))=\bigcap_t f(\mathrm D(t)),
\]
et par hypothèse, pour tout $t\in B-\mathfrak q$, $y$ est intérieur à
$f(\mathrm D(t))$, donc $f(\mathrm D(t))$ contient
$\operatorname{Spec}(\mathcal O_{Y,y})$~; d'où
$\operatorname{Spec}(\mathcal O_{Y,y})
\subset f(\operatorname{Spec}(\mathcal O_{X,x}))$, ce qui établit
notre assertion. Enfin, lorsque $f$ est localement de présentation finie, b)
implique a) en vertu de (1.10.2) appliqué à la restriction $U\to Y$ de $f$ à
un voisinage ouvert arbitraire $U$ de $x$.
\end{proof}

\begin{corollary}[1.10.4]
\label{IV.1.10.4-fr}
Soit $f:X\to Y$ un morphisme. Considérons les conditions suivantes~:
\begin{enumerate}
  \item[a)] $f$ est ouvert.
  \item[b)] Pour tout $x\in X$ et toute générisation $y'$ de $y=f(x)$, il
  existe une générisation $x'$ de $x$ telle que $f(x')=y'$.
  \item[b')] Pour tout $x\in X$, l'image par $f$ de
  $\operatorname{Spec}(\mathcal O_{X,x})$ est
  $\operatorname{Spec}(\mathcal O_{Y,f(x)})$.
  \item[c)] Pour toute partie fermée irréductible $Y'$ de $Y$, toute
  composante irréductible de $f^{-1}(Y')$ domine $Y'$.
\end{enumerate}
Alors on a les implications
$\mathrm{a)}\Rightarrow\mathrm{b)}\Leftrightarrow\mathrm{b')}
\Leftrightarrow\mathrm{c)}$. Si en outre $f$ est localement de présentation
finie, les quatre conditions sont équivalentes.
\end{corollary}

\begin{proof}
Dire que $f$ est ouvert signifie que $f$ est ouvert en tout point $x\in X$,
donc les implications
$\mathrm{a)}\Rightarrow\mathrm{b)}\Leftrightarrow\mathrm{b')}$ résultent des
implications analogues dans (1.10.3), ainsi que l'implication
$\mathrm{b)}\Rightarrow\mathrm{a)}$ lorsque $f$ est localement de présentation
finie. L'implication $\mathrm{c)}\Rightarrow\mathrm{b)}$ résulte aussi de
l'implication analogue dans (1.10.3)~; prouvons enfin que b) entraîne c). En
effet, soit $x'$ le point générique d'une composante
irréductible de $f^{-1}(Y')$, et montrons que $y'=f(x')$ est le point générique
de $Y'$. Soit $y''$ une générisation de $y'$~; il existe par hypothèse une
générisation $x''$ de $x'$ telle que $f(x'')=y''$, et comme
$x''\in f^{-1}(Y')$, on a nécessairement $x''=x'$, donc $y''=y'$, ce qui
achève la démonstration.
\end{proof}

\begin{center}
\emph{(A suivre.)}
\end{center}

\clearpage
\oldpage[IV-1]{251}
\phantomsection
\addcontentsline{toc}{section}{Bibliographie (suite)}
\section*{BIBLIOGRAPHIE (\emph{suite})}

\begingroup
\small
\setlength{\parindent}{0pt}
\setlength{\parskip}{3pt}
[30] M.~Nagata, \emph{Local rings}, \emph{Interscience Tracts}, 13 (1962),
Interscience, New York.

[31] M.~Nagata, \emph{A Jacobian criterion of simple points},
\emph{Ill. Journ. of Math.}, 1 (1957), p.~427--432.

[32] A.~Néron, \emph{Modèles minimaux des variétés abéliennes sur les corps
locaux et globaux}, \emph{Publ. Inst. Hautes Études Sci.}, n\textsuperscript{o}
21 (1964).

\clearpage
\oldpage[IV-1]{252}
\phantomsection
\addcontentsline{toc}{section}{Index des notations}
\section*{INDEX DES NOTATIONS}

$\dim_c(X)$, $\dim(X)$ ($X$ espace topologique) : 0, 14.1.2.\par
$\dim_x(X)$ ($X$ espace topologique, $x$ point de $X$) : 0, 14.1.2.\par
$\operatorname{codim}(Y,X)$ ($X$ espace topologique, $Y$ partie fermée de $X$) : 0, 14.2.1.\par
$\operatorname{codim}_x(Y,X)$ ($X$ espace topologique, $Y$ partie fermée de $X$, $x$ point de $X$) : 0, 14.2.4.\par
$N[T_1,\ldots,T_n]$ ($N$ $A$-module) : 0, 15.1.1.\par
$\dim(A)$ ($A$ anneau) : 0, 16.1.1.\par
$\operatorname{ht}(\mathfrak a)$ ($\mathfrak a$ idéal) : 0, 16.1.3.\par
$\dim_A(M)$, $\dim(M)$ ($M$ $A$-module) : 0, 16.1.7.\par
$P_{\mathfrak q}(M,n)$ ($M$ module de type fini sur un anneau semi-local noethérien $A$, $\mathfrak q$ idéal de définition de $A$) : 0, 16.2.1.\par
$\operatorname{prof}_A(M)$, $\operatorname{prof}(M)$ ($M$ $A$-module) : 0, 16.4.5.\par
$\operatorname{coprof}_A(M)$, $\operatorname{coprof}(M)$ ($M$ $A$-module) : 0, 16.4.9.\par
$\operatorname{dim.proj}_A(M)$, $\operatorname{dim.proj}(M)$,
$\operatorname{dim.inj}_A(M)$, $\operatorname{dim.inj}(M)$ ($M$ $A$-module) : 0, 17.2.1.\par
$\operatorname{dim.coh}(A)$ ($A$ anneau) : 0, 17.2.8.\par
$\operatorname{dim.proj}(\mathscr F)$, $\operatorname{dim.inj}(\mathscr F)$
($\mathscr F$ Module sur un espace annelé) : 0, 17.2.14.\par
$\operatorname{dim.coh}(X)$ ($X$ espace annelé) : 0, 17.2.14.\par
$E\times_B F$ ($B$, $E$, $F$ $A$-anneaux) : 0, 18.1.2.\par
$D_B(L)$ ($B$ anneau, $L$ $B$-bimodule) : 0, 18.2.3.\par
$Y\oplus_X Z$ ($X$, $Y$, $Z$ $A$-bimodules) : 0, 18.2.7.\par
$\operatorname{Exan}_A(B,L)$ ($B$ $A$-anneau, $L$ $B$-bimodule) : 0, 18.3.4.\par
$\operatorname{Exan}_{A/A'}(B,L)$ : 0, 18.3.7.\par
$\operatorname{Exal}_A(B,L)$, $\operatorname{Exal}_{A/A'}(B,L)$ ($A$ anneau commutatif, $B$ $A$-algèbre, $L$ $B$-bimodule) : 0, 18.4.1.\par
$\operatorname{Exalcom}_A(B,L)$, $\operatorname{Exalcom}_{A/A'}(B,L)$ ($A$ anneau commutatif, $B$ $A$-algèbre commutative, $L$ $B$-module) : 0, 18.4.2.\par
$\mathrm H_A^2(B,L)$, $\mathrm H_A^2(B,L)^s$ : 0, 18.4.3.\par
$\operatorname{Exantop}_A(B,L)$, $\operatorname{Exaltop}_A(B,L)$,
$\operatorname{Exalcotop}_A(B,L)$ : 0, 18.5.1.\par
$\operatorname{Exantop}_{A/A'}(B,L)$, $\operatorname{Exaltop}_{A/A'}(B,L)$,
$\operatorname{Exalcotop}_{A/A'}(B,L)$ : 0, 18.5.2.\par
$\operatorname{Hom.\ cont}_C(M,N)$ ($M$, $N$ $C$-bimodules topologiques) : 0, 18.5.3.\par
$\operatorname{Dér}_A(B,L)$ ($B$ $A$-anneau, $L$ $B$-bimodule) : 0, 20.1.2.\par
$\operatorname{Dér.\ cont}_A(B,L)$ : 0, 20.3.1.\par
$p_{B/A}$, $\mathfrak J_{B/A}$ ($B$ $A$-algèbre topologique commutative) : 0, 20.4.1.\par
$P^1_{B/A}$, $\varepsilon_{B/A}$ ($B$ $A$-algèbre topologique commutative) : 0, 20.4.2.\par
$\Omega^1_{B/A}$, $\widehat{\Omega}^1_{B/A}$ ($B$ $A$-algèbre topologique commutative) : 0, 20.4.3.\par
$\Omega^1_B$ ($B$ anneau topologique) : 0, 20.4.3.\par
$d_{B/A}$, $d(x)$, $dx$, $d_B$ : 0, 20.4.6.\par
$\Upsilon_{C/B/A}$, $\Upsilon_{C/B}$ ($A$, $B$, $C$ anneaux) : 0, 20.6.1.\par
$\chi_E$ ($E$ $B$-extension $A$-triviale de $C$ par $L$) : 0, 20.6.8.\par

\clearpage
\oldpage[IV-1]{253}
$\Upsilon^C_{B/A/\Lambda}$, $\Upsilon^C_{B/A}$ ($\Lambda$, $A$, $B$, $C$ anneaux) : 0, 20.6.14.\par
$\chi_{B/A}$, $\chi_B$ ($B$ $A$-algèbre) : 0, 20.6.24.\par
$\chi_E$ ($E$ $B$-algèbre topologique) : 0, 20.7.10.\par
$F_A$, $A^p$, $A^{(p)}$, $E^{(p)}$ ($A$ anneau de caractéristique $p>0$, $E$ $A$-module) : 0, 21.1.4.\par
$\pi_{B/A}$, $\Theta_{B/A}$, $\Xi_{B/A}$ ($A$, $B$ anneaux de caractéristique $p>0$) : 0, 21.3.2.\par
$\Delta(L/K,k/k_0)$, $d(L/K,k/k_0)$, $\Delta(L/K,k)$, $d(L/K,k)$ ($k_0$, $k$, $K$, $L$ corps) : 0, 21.6.1.\par
$\chi'_{B/A}$ ($A$, $B$ anneaux de caractéristique $p>0$) : 0, 22.4.6.\par
$\mathscr T^{\mathrm{cons}}$, $X^{\mathrm{cons}}$ ($X$ préschéma, $\mathscr T$ topologie de $X$) : IV, 1.9.12.\par
$f^{\mathrm{cons}}$ ($f$ morphisme de préschémas) : IV, 1.9.13.\par

\clearpage
\oldpage[IV-1]{254}
\phantomsection
\addcontentsline{toc}{section}{Index terminologique}
\section*{INDEX TERMINOLOGIQUE}

Algèbre augmentée des parties principales d'ordre 1 : 0, 20.4.2.\par
$A$-algèbre de Cohen : 0, 19.8.1.\par
Algèbre de présentation finie : IV, 1.4.1.\par
$A$-algèbre essentiellement de type fini : IV, 1.3.8.\par
$A$-anneau : 0, 18.1.1.\par
$A$-anneau augmenté sur $B$ : 0, 18.1.4.\par
Anneau biéquidimensionnel : 0, 16.1.4.\par
Anneau caténaire : 0, 16.1.4.\par
Anneau de caractéristique $p$ : 0, 21.1.1.\par
Anneau de Cohen : 0, 19.8.4.\par
Anneau de Cohen-Macaulay : 0, 16.5.1 et 16.5.13.\par
Anneau équicodimensionnel, équidimensionnel : 0, 16.1.4.\par
Anneau local géométriquement unibranche : 0, 23.1.7.\par
Anneau japonais : 0, 23.1.1.\par
Anneau local premier, anneau local complet premier : 0, 19.8.3.\par
Anneau régulier : 0, 17.1.1 et 17.3.6.\par
Anneau local unibranche : 0, 23.1.7.\par
Anneau universellement japonais : 0, 23.1.1.\par
Application ouverte en un point : IV, 1.10.2.\par
Augmentation : 0, 18.1.4.\par

Bimorphisme formel : 0, 19.1.2.\par

Chaîne dans un ensemble ordonné, chaîne de parties fermées irréductibles d'un espace topologique : 0, 14.1.1.\par
Chaîne saturée : 0, 14.3.1.\par
Codimension combinatoire, codimension d'une partie fermée d'un espace topologique : 0, 14.2.1.\par
Codimension en un point d'une partie fermée d'un espace topologique : 0, 14.2.4.\par
Condition des chaînes : 0, 14.3.2.\par
Coprofondeur d'un module sur un anneau local noethérien : 0, 16.4.9.\par
Corps $k_0$-admissible, corps admissible pour une extension : 0, 21.6.1.\par
Corps de multiplicité radicielle finie sur un sous-corps : 0, 19.6.6.\par

Défaut de $k_0$-admissibilité d'une extension : 0, 21.6.1.\par
Dérivation, $A$-dérivation : 0, 20.1.2.\par
Différentielles, 1-différentielles, différentielles absolues : 0, 20.4.3.\par
Différentielle extérieure : 0, 20.4.6.\par
Dimension combinatoire, dimension d'un espace topologique : 0, 14.1.2.\par
Dimension d'un espace topologique en un point : 0, 14.1.2.\par
Dimension de Krull, dimension d'un anneau : 0, 16.1.1.\par
Dimension d'un module : 0, 16.1.7.\par
Dimension cohomologique, dimension cohomologique globale d'un anneau : 0, 17.2.8.\par
Dimension cohomologique d'un espace annelé en un point, dimension cohomologique d'un espace annelé : 0, 17.2.14.\par

\clearpage
\oldpage[IV-1]{255}
Dimension injective, dimension projective d'un module : 0, 17.2.1.\par
Dimension injective en un point, dimension injective d'un faisceau : 0, 17.2.14.\par
Dimension projective en un point, dimension projective d'un faisceau : 0, 17.2.14.\par

Élément de $A$ non diviseur de 0 dans un $A$-module : 0, 15.1.1.\par
Élément $M$-régulier, élément $M$-quasi-régulier : 0, 15.1.4.\par
Ensemble ind-constructible, ensemble pro-constructible : IV, 1.9.4.\par
Épimorphisme formel : 0, 19.1.2.\par
$A$-équivalence de $A$-extensions : 0, 18.2.2.\par
Espace biéquidimensionnel : 0, 14.3.3.\par
Espace caténaire : 0, 14.3.2.\par
Espace équicodimensionnel : 0, 14.2.1.\par
Espace équidimensionnel : 0, 14.1.3.\par
$A$-extension d'un $A$-anneau $B$ par un $B$-bimodule : 0, 18.2.2.\par
Extension déduite d'une autre par un homomorphisme de bimodules : 0, 18.2.8.\par
Extension de Hochschild : 0, 18.4.3.\par
$B$-extension $A$-triviale : 0, 18.3.7.\par
Extension triviale type : 0, 18.2.3.\par
$A$-extensions $A$-équivalentes : 0, 18.2.2.\par

Famille $p$-libre, famille absolument $p$-libre : 0, 21.1.9.\par
Formellement étale (algèbre) : 0, 19.10.2.\par
Formellement inversible à droite (homomorphisme) : 0, 19.1.15.\par
Formellement inversible à gauche (homomorphisme) : 0, 19.1.5.\par
Formellement non ramifiée (algèbre) : 0, 19.10.2.\par
Formellement projectif (module) : 0, 19.2.1.\par
Formellement lisse (algèbre) : 0, 19.3.1.\par
Formellement lisse relativement à un anneau (algèbre) : 0, 19.9.1.\par

Géométriquement régulier sur un corps (anneau local) : 0, 19.6.5.\par
Groupe des classes de $A$-extensions : 0, 18.3.4.\par

Hauteur d'un idéal : 0, 16.1.3.\par
$A$-homomorphisme de $A$-anneaux : 0, 18.1.1.\par
Homomorphisme caractéristique d'une $B$-extension $A$-triviale : 0, 20.6.8.\par
Homomorphisme caractéristique d'une $A$-algèbre relativement à un anneau et à un idéal : 0, 20.6.24.\par
Homomorphisme caractéristique d'une algèbre de caractéristique $p>0$ : 0, 22.4.6.\par

Idéal d'augmentation : 0, 18.1.4.\par
Image réciproque d'un anneau augmenté : 0, 18.1.5.\par
Image réciproque d'une $A$-extension : 0, 18.2.5.\par

Longueur d'une chaîne : 0, 14.1.1.\par

Module biéquidimensionnel, caténaire, équicodimensionnel, équidimensionnel : 0, 16.1.7.\par
Module de Cohen-Macaulay : 0, 16.5.1 et 16.5.13.\par
Module d'imperfection : 0, 20.6.1.\par
Monomorphisme formel : 0, 19.1.2.\par
Morphisme de $A$-extensions : 0, 18.2.4.\par
Morphisme de présentation finie en un point, morphisme localement de présentation finie : IV, 1.4.2.\par
Morphisme de présentation finie : IV, 1.6.1.\par
Morphisme quasi-séparé : IV, 1.2.1.\par

$p$-anneau de Cohen : 0, 19.8.4.\par
$p$-base, $p$-base absolue : 0, 21.1.9.\par
Point maximal d'un préschéma : IV, 1.1.4.\par

\clearpage
\oldpage[IV-1]{256}
Polynôme de Hilbert-Samuel : 0, 16.2.1.\par
Préschéma de présentation finie sur un autre : IV, 1.6.1.\par
Préschéma quasi-séparé : IV, 1.2.1.\par
Produit fibré de $A$-anneaux : 0, 18.1.2.\par
Profondeur d'un module : 0, 16.4.5.\par

Somme amalgamée de bimodules : 0, 18.2.7.\par
Strictement formellement projectif (module) : 0, 19.2.3.\par
Suite $M$-régulière, suite $M$-quasi-régulière : 0, 15.1.7.\par
Suite $\mathscr F$-régulière, suite $\mathscr F$-quasi-régulière : 0, 15.2.2.\par
Système de paramètres pour un module sur un anneau semi-local noethérien : 0, 16.3.6.\par
Système de $p$-générateurs, système absolu de $p$-générateurs : 0, 21.1.9.\par
Système fondamental d'idéaux ouverts, de sous-modules ouverts : 0, 19.0.2.\par
Système régulier de paramètres d'un anneau local régulier : 0, 17.1.6.\par

Topologie constructible sur un préschéma : IV, 1.9.12.\par
Topologie sur un $A$-module déduite de la topologie de $A$ : 0, 19.0.2.\par
Trivial ($A$-anneau augmenté) : 0, 18.1.4.\par
$A$-triviale (extension) : 0, 18.2.3 et 18.3.7.\par

\clearpage
\oldpage[IV-1]{257}
\phantomsection
\addcontentsline{toc}{section}{Table des matières}
\section*{TABLE DES MATIÈRES}

\newcommand{\ivonecontentsline}[3]{\noindent\makebox[17mm][l]{#1}\parbox[t]{0.77\linewidth}{#2\dotfill}\makebox[12mm][r]{#3}\par}
\ivonecontentsline{Chapitre 0.}{\textbf{Préliminaires} (\emph{suite})}{5}
\ivonecontentsline{\S~14.}{Dimension combinatoire d'un espace topologique}{6}
\ivonecontentsline{14.1.}{Dimension combinatoire d'un espace topologique}{6}
\ivonecontentsline{14.2.}{Codimension d'une partie fermée}{8}
\ivonecontentsline{14.3.}{La condition des chaînes}{10}
\ivonecontentsline{\S~15.}{Suites $M$-régulières et suites $\mathscr F$-régulières}{12}
\ivonecontentsline{15.1.}{Suites $M$-régulières et suites $M$-quasi-régulières}{12}
\ivonecontentsline{15.2.}{Suites $\mathscr F$-régulières}{20}
\ivonecontentsline{\S~16.}{Dimension et profondeur dans les anneaux locaux noethériens}{22}
\ivonecontentsline{16.1.}{Dimension d'un anneau}{22}
\ivonecontentsline{16.2.}{Dimension d'un anneau semi-local noethérien}{25}
\ivonecontentsline{16.3.}{Systèmes de paramètres dans un anneau local noethérien}{28}
\ivonecontentsline{16.4.}{Profondeur et coprofondeur}{32}
\ivonecontentsline{16.5.}{Modules de Cohen-Macaulay}{36}
\ivonecontentsline{\S~17.}{Anneaux réguliers}{39}
\ivonecontentsline{17.1.}{Définition des anneaux réguliers}{39}
\ivonecontentsline{17.2.}{Rappels sur la dimension projective et la dimension injective des modules}{42}
\ivonecontentsline{17.3.}{Théorie cohomologique des anneaux réguliers}{46}
\ivonecontentsline{\S~18.}{Compléments sur les extensions d'algèbres}{51}
\ivonecontentsline{18.1.}{Images réciproques d'anneaux augmentés}{51}
\ivonecontentsline{18.2.}{Extensions d'un anneau par un bimodule}{54}
\ivonecontentsline{18.3.}{Le groupe des classes de $A$-extensions}{59}
\ivonecontentsline{18.4.}{Extensions d'algèbres}{64}
\ivonecontentsline{18.5.}{Cas des anneaux topologiques}{66}
\ivonecontentsline{\S~19.}{Algèbres formellement lisses et anneaux de Cohen}{69}
\ivonecontentsline{19.0.}{Introduction}{69}
\ivonecontentsline{19.1.}{Épimorphismes et monomorphismes formels}{71}
\ivonecontentsline{19.2.}{Modules formellement projectifs}{78}

\clearpage
\oldpage[IV-1]{258}
\ivonecontentsline{19.3.}{Algèbres formellement lisses}{79}
\ivonecontentsline{19.4.}{Premiers critères de lissité formelle}{86}
\ivonecontentsline{19.5.}{Lissité formelle et anneaux gradués associés}{90}
\ivonecontentsline{19.6.}{Cas des algèbres sur un corps}{100}
\ivonecontentsline{19.7.}{Cas des homomorphismes locaux : théorèmes d'existence et d'unicité}{104}
\ivonecontentsline{19.8.}{Algèbres de Cohen et $p$-anneaux de Cohen; application à la structure des anneaux locaux complets}{109}
\ivonecontentsline{19.9.}{Algèbres relativement formellement lisses}{114}
\ivonecontentsline{19.10.}{Algèbres formellement non ramifiées et algèbres formellement étales}{115}
\ivonecontentsline{\S~20.}{Dérivations et différentielles}{116}
\ivonecontentsline{20.1.}{Dérivations et extensions d'algèbres}{117}
\ivonecontentsline{20.2.}{Propriétés fonctorielles des dérivations}{119}
\ivonecontentsline{20.3.}{Dérivations continues dans les anneaux topologiques}{121}
\ivonecontentsline{20.4.}{Parties principales et différentielles}{123}
\ivonecontentsline{20.5.}{Propriétés fonctorielles fondamentales de $\Omega^1_{B/A}$}{128}
\ivonecontentsline{20.6.}{Modules d'imperfection et homomorphismes caractéristiques}{136}
\ivonecontentsline{20.7.}{Généralisations aux anneaux topologiques}{147}
\ivonecontentsline{\S~21.}{Différentielles dans les anneaux de caractéristique $p$}{153}
\ivonecontentsline{21.1.}{Systèmes de $p$-générateurs et $p$-bases}{154}
\ivonecontentsline{21.2.}{$p$-bases et lissité formelle}{157}
\ivonecontentsline{21.3.}{$p$-bases et modules d'imperfection}{160}
\ivonecontentsline{21.4.}{Cas des extensions de corps}{162}
\ivonecontentsline{21.5.}{Application : critères de séparabilité}{164}
\ivonecontentsline{21.6.}{Corps admissibles pour une extension}{167}
\ivonecontentsline{21.7.}{L'égalité de Cartier}{169}
\ivonecontentsline{21.8.}{Critères d'admissibilité}{171}
\ivonecontentsline{21.9.}{Modules de différentielles complétés dans les anneaux de séries formelles}{176}
\ivonecontentsline{\S~22.}{Critères différentiels de lissité formelle et de régularité}{182}
\ivonecontentsline{22.1.}{Relèvement de la lissité formelle}{183}
\ivonecontentsline{22.2.}{Caractérisation différentielle des algèbres locales formellement lisses sur un corps}{186}
\ivonecontentsline{22.3.}{Application aux relations entre certains anneaux locaux et leurs complétés}{191}
\ivonecontentsline{22.4.}{Résultats préliminaires sur les extensions finies d'anneaux locaux dont l'idéal maximal est de carré nul}{193}

\clearpage
\oldpage[IV-1]{259}
\ivonecontentsline{22.5.}{Algèbres géométriquement régulières et algèbres formellement lisses}{201}
\ivonecontentsline{22.6.}{Critère jacobien de Zariski}{205}
\ivonecontentsline{22.7.}{Le critère jacobien de Nagata}{209}
\ivonecontentsline{\S~23.}{Anneaux japonais}{213}
\ivonecontentsline{23.1.}{Anneaux japonais}{213}
\ivonecontentsline{23.2.}{Clôture intégrale d'un anneau local noethérien intègre}{217}
\ivonecontentsline{}{\textsc{Chapitre IV.} --- \textbf{Étude locale des schémas et des morphismes de schémas}}{222}
\ivonecontentsline{\S~1.}{Conditions de finitude relatives. Ensembles constructibles dans les préschémas}{224}
\ivonecontentsline{1.1.}{Morphismes quasi-compacts}{224}
\ivonecontentsline{1.2.}{Morphismes quasi-séparés}{226}
\ivonecontentsline{1.3.}{Morphismes localement de type fini}{228}
\ivonecontentsline{1.4.}{Morphismes localement de présentation finie}{230}
\ivonecontentsline{1.5.}{Morphismes de type fini}{233}
\ivonecontentsline{1.6.}{Morphismes de présentation finie}{234}
\ivonecontentsline{1.7.}{Amélioration de résultats antérieurs}{236}
\ivonecontentsline{1.8.}{Morphismes de présentation finie et ensembles constructibles}{238}
\ivonecontentsline{1.9.}{Ensembles pro-constructibles et ind-constructibles}{241}
\ivonecontentsline{1.10.}{Applications aux morphismes ouverts}{249}
\ivonecontentsline{}{Bibliographie (\emph{suite})}{251}
\ivonecontentsline{}{Index des notations}{252}
\ivonecontentsline{}{Index terminologique}{254}

\vspace{2em}
\begin{flushright}
\emph{Reçu le 15 mars 1963.}
\end{flushright}
\let\ivonecontentsline\relax
\endgroup
